\chapter{图像的特征与分割}
\label{chap:vision-image-features}

设想一辆车在路口拍下一幅街景：车道线斜穿画面，路边有一块招牌，人行道上站着两位行人，上半部是天空，下半部是路面。我们希望计算机回答几个具体问题：车道线在哪里、朝哪个方向；同一块招牌在另一幅从不同位置、不同距离拍的照片中位于何处；路面与天空各占哪些像素；画面中有没有行人、行人在哪里。数组$\bm X\in\mathbb R^{h\times w\times c}$中并没有“车道线”“招牌”或“行人”这样的条目，只有每个位置记录的亮度与颜色。要回答这些问题，必须先从像素值的空间变化中提取出比单个像素更稳定、更有区分力的量，再把这些量组织成直线、对应点、区域或检测框。

这一章研究在不借助大规模学习的条件下怎样完成这件事。读者需要的前置知识包括\ref{chap:vision-image-transforms}章的线性滤波、梯度算子、图像金字塔与形态学运算，以及\ref{chap:vision-image-formation}章的像素邻接与坐标约定。本章所用的数学工具大多在卷一、卷二与卷四出现过，例如卷积、对称矩阵的特征分解、图上的最小割与谱松弛、K均值聚类与支持向量机；每次用到时先在本章写出所需结论和条件，再给出参考位置。卷一、卷二尚未覆盖的工具，例如结构张量、高斯尺度空间和图像矩，在本章用够用的篇幅讲清楚。

全章的线索是像素值的空间变化。灰度沿某个方向突变的地方是边缘，边缘连成线就给出车道线的位置；在两个方向上都变化的地方是角点或斑点，它们能在不同照片中被重复找到，适合做匹配；局部梯度方向的统计分布描述了形状与纹理，可以用来比较两个窗口是否属于同一类对象。区域是另一种组织方式：把灰度、颜色或纹理相近而又彼此相连的像素归为一组，得到路面、天空等区域，再用面积、边界和矩等数值描述区域的形状。最后，把局部特征与分类器结合，就得到不依赖深度网络的检测与识别系统。这些经典方法的设计选择（用什么算子、在多大尺度上统计、怎样归一化、怎样在全图搜索）后来大多以可学习的形式出现在卷积网络和视觉Transformer中；理解它们为什么这样设计、在什么条件下失效，有助于理解学习方法改变了什么、又保留了什么。

本章约定：灰度图像记作定义在平面上的函数$I(x,y)$，$x$向右、$y$向下；离散时像素位置为整数。图像梯度记作$\nabla I=(I_x,I_y)^{\top}$，$I_x=\partial I/\partial x$。高斯函数
\begin{equation}
  G_{\sigma}(x,y)=\frac{1}{2\pi\sigma^2}\exp\!\Bigl(-\frac{x^2+y^2}{2\sigma^2}\Bigr)
  \label{eq:vision-features-gaussian}
\end{equation}
的标准差$\sigma$称为尺度，$*$表示二维卷积。卷积满足交换律与结合律，并与求导可交换：$\partial(G_{\sigma}*I)/\partial x=(\partial G_{\sigma}/\partial x)*I$。这一性质在连续可积函数和离散有限支撑核上都成立，它使“先平滑再求导”可以合并为一次与导数核的卷积，可参考\BookChapterRef{chap:signal-analysis}{卷一“信号”章}。

\section{边缘与直线的检测}
\label{sec:vision-features-edges}

街景中车道线与路面、车身与背景、招牌与墙面之间都有清晰的分界。分界两侧的像素各自比较均匀，跨过分界时灰度在很短距离内发生大幅变化。\term{边缘}（edge）指图像中灰度沿某一方向急剧变化的位置，\term{边缘检测}（edge detection）就是找出这些位置，通常还给出变化的强度与方向。输入是一幅灰度图像，输出是一张二值或带强度的边缘图；若要得到车道线这样的几何对象，还需要把零散的边缘点组织成直线或曲线。

难点来自三个方面。第一，真实边缘并不是理想的阶跃：镜头离焦和像素积分使边缘变成宽度不一的斜坡，细线（如车道线本身）则在两侧都有灰度跳变。第二，求导会放大噪声，因为噪声在相邻像素之间独立变化，而差分恰好突出这种快速变化。第三，边缘点本身是零散的，阴影、纹理和遮挡会使同一条直线断成若干段，又混入许多不属于直线的点。

围绕这些难点有四条路线。第一条用一阶导数：边缘处梯度幅值最大，对梯度幅值取阈值即得边缘。第二条用二阶导数：一阶导数的极大处正是二阶导数的零交叉，零交叉的位置不依赖阈值，配合高斯平滑还可以选择观察的尺度。第三条是Canny检测器，它把检测、定位与单一响应写成三个可计算的准则，由此推出近似最优的平滑求导算子，再用非极大值抑制细化边缘、用双阈值连接断开的边缘。前三条路线回答“哪些像素是边缘”；第四条路线Hough变换回答“哪些边缘点属于同一条直线或圆”，它让每个边缘点为所有可能经过它的直线投票，用投票的累积来对抗断裂与干扰。

\subsection{基于一阶梯度的检测}
\label{sec:vision-features-gradient-edges}

一阶路线的出发点是：沿垂直于边缘的方向看，灰度剖面像一个斜坡，斜坡中点处一阶导数的绝对值最大。二维图像中，函数$I$在点$(x,y)$处沿单位方向$\bm u=(\cos\alpha,\sin\alpha)^{\top}$的方向导数为$\nabla I^{\top}\bm u$。由柯西—施瓦茨不等式，$\lvert\nabla I^{\top}\bm u\rvert\leq\lVert\nabla I\rVert_2$，等号在$\bm u$与$\nabla I$同向或反向时取得。因此梯度方向就是灰度变化最快的方向，梯度幅值就是该方向上的变化率：
\begin{equation}
  M(x,y)=\sqrt{I_x^2+I_y^2},\qquad
  \alpha(x,y)=\operatorname{atan2}(I_y,I_x).
  \label{eq:vision-features-gradient-magnitude}
\end{equation}
边缘本身沿着与梯度垂直的方向延伸。车道线的一侧边缘若从左下斜向右上，那么该处的梯度就指向跨越车道线边界的方向。

离散图像上用差分近似偏导数。中心差分$I_x(x,y)\approx[I(x+1,y)-I(x-1,y)]/2$对应卷积核$\tfrac12(-1,0,1)$；Sobel算子把这个差分核与垂直方向的平滑核$(1,2,1)^{\top}$相乘，得到可分离的$3\times3$核
\[
  \bm K_x=\begin{pmatrix}-1&0&1\\-2&0&2\\-1&0&1\end{pmatrix},\qquad
  \bm K_y=\bm K_x^{\top}.
\]
这两种算子的构造与可分离性已在\ref{sec:vision-transforms-sharpening}节介绍，这里关心的是它们在噪声下检测边缘的能力。

先看一个算例。设$3\times3$邻域为
\[
  \begin{pmatrix}10&10&10\\10&10&50\\10&50&50\end{pmatrix},
\]
它表示一条从右上斜向左下的亮暗分界，右下方较亮。用Sobel核做互相关（核不翻转，这是图像处理中的通常做法）：$\bm K_x$的三行分别给出$-10+10=0$、$-20+100=80$、$-10+50=40$，所以$I_x=120$；$\bm K_y$的上下两行分别给出$-10-20-10=-40$与$10+100+50=160$，所以$I_y=120$。于是$M=120\sqrt2\approx169.7$，$\alpha=45^\circ$。由于$y$轴向下，$45^\circ$表示梯度指向右下方，即由暗区指向亮区，边缘沿右上—左下方向延伸，与邻域中的分界一致。

为什么Sobel比中心差分更适合检测边缘？可以比较二者的信噪比。设每个像素的观测是真实灰度加上独立同分布、方差为$s^2$的噪声；线性滤波输出的噪声方差等于核系数平方和乘以$s^2$，这一结论只用到独立随机变量之和的方差可加。不带系数$\tfrac12$的中心差分$(-1,0,1)$的系数平方和为$2$；Sobel核的系数平方和为$(1^2+2^2+1^2)(1^2+0+1^2)=12$。对于跨越垂直分界、高度为$h$的理想阶跃，中心差分的响应为$h$，Sobel的三行各给出$h,2h,h$，合计$4h$。二者的信噪比分别为
\[
  \frac{h}{\sqrt2\,s}\approx0.71\,\frac hs,\qquad
  \frac{4h}{\sqrt{12}\,s}\approx1.15\,\frac hs.
\]
Sobel沿边缘方向平均三行，把信号相干地累加，而噪声只按平方和累加，信噪比因此提高。代价是边缘在平滑方向上变宽，相邻边缘更容易混在一起。这正是后面所有边缘检测器都要面对的取舍：平滑越强，噪声越小，定位越粗。

得到$M$之后，最简单的做法是取阈值$T$，令$M(x,y)\geq T$的像素为边缘。这样做有两个直接问题。其一，阶跃经平滑后变成斜坡，斜坡上的若干像素梯度都超过阈值，得到的边缘有好几个像素宽，无法精确说明边缘在哪里。其二，阈值难以兼顾：取高了，对比度较低的边缘（如阴影中的车道线）会断开；取低了，路面纹理和噪声会产生大量零碎响应。只依靠单个像素的梯度幅值，无法区分“处在一条连续轮廓上的弱边缘”和“孤立的噪声点”。这两个问题分别由后面的非极大值抑制和滞后阈值回答。

\subsection{基于二阶导数零交叉的检测}
\label{sec:vision-features-log}

一阶路线需要先定阈值再找边缘，而边缘的位置本应与阈值无关。换一个角度看，斜坡边缘的位置是一阶导数取极大的点，在那里二阶导数由正变负或由负变正，即发生\term{零交叉}（zero crossing）。只要找出二阶导数的符号变化，就能在不规定幅值阈值的情况下确定边缘位置。

二维中通常使用各向同性的拉普拉斯算子$\nabla^2I=I_{xx}+I_{yy}$。它与方向无关，一次计算就能覆盖所有方向的边缘。二阶差分对噪声比一阶差分更敏感，所以必须先平滑。Marr与Hildreth提出先用高斯核平滑，再求拉普拉斯\scite{vision-marr1980}。由于卷积与求导可交换，
\[
  \nabla^2(G_{\sigma}*I)=(\nabla^2G_{\sigma})*I,
\]
平滑与求二阶导可以合并为与一个核的卷积。这个核称为\term{高斯拉普拉斯}（Laplacian of Gaussian, LoG）。它的表达式由逐项求导得到：对式\eqref{eq:vision-features-gaussian}，$\partial G_{\sigma}/\partial x=-(x/\sigma^2)G_{\sigma}$，再求一次导，由乘积法则得$\partial^2G_{\sigma}/\partial x^2=(x^2/\sigma^4-1/\sigma^2)G_{\sigma}$；$y$方向同理。两式相加，
\begin{equation}
  \nabla^2G_{\sigma}(x,y)=\frac{x^2+y^2-2\sigma^2}{\sigma^4}\,G_{\sigma}(x,y).
  \label{eq:vision-features-log}
\end{equation}
这个核在半径$r=\sqrt{x^2+y^2}<\sqrt2\sigma$的圆盘内为负，圆盘外为正，形如倒置的“墨西哥草帽”。

一维算例可以看出零交叉怎样确定位置。设一行像素经过平滑后的灰度为$(10,10,12,20,40,48,50,50)$。二阶差分$f_{i-1}-2f_i+f_{i+1}$在第二至第七个位置依次为$2,6,12,-12,-6,-2$，符号在第四、五个位置之间由正变负，所以边缘位于灰度$20$与$40$之间。一阶差分$f_{i+1}-f_i$为$0,2,8,20,8,2,0$，最大值$20$也正好落在这两个像素之间。零交叉不需要阈值就给出了这个位置；实际中仍会要求零交叉两侧的差值超过一个小阈值，以免平坦区域中噪声引起的微小符号变化也被当成边缘。

尺度$\sigma$决定了在多粗的分辨率上观察边缘。$\sigma$越大，越多的细小纹理和噪声被平均掉，只剩下大尺度的轮廓；但相距小于约$\sigma$的两条边缘会被平滑成一个斜坡，拐角也被抹圆。没有一个$\sigma$对所有结构都合适：车道线的宽度在近处有几十个像素，在远处只有一两个像素。Marr与Hildreth因此主张在多个尺度上检测边缘再综合结果。这个“在多个尺度上观察同一幅图像”的想法，在\ref{sec:vision-features-sift}节发展为尺度空间。

LoG与高斯核的尺度导数之间有一个后面要用到的关系。对式\eqref{eq:vision-features-gaussian}按$\sigma$求偏导，
\[
  \frac{\partial G_{\sigma}}{\partial\sigma}
  =\Bigl(-\frac{2}{\sigma}+\frac{x^2+y^2}{\sigma^3}\Bigr)G_{\sigma}
  =\sigma\,\nabla^2G_{\sigma},
\]
第二个等号只需把式\eqref{eq:vision-features-log}乘以$\sigma$比较。用差商近似导数，对$k>1$有$G_{k\sigma}-G_{\sigma}\approx(k\sigma-\sigma)\,\partial G_{\sigma}/\partial\sigma=(k-1)\sigma^2\nabla^2G_{\sigma}$。也就是说，两个尺度相近的高斯核之差近似于乘了$\sigma^2$的LoG。这个\term{高斯差分}（difference of Gaussians, DoG）只需两次高斯平滑和一次相减，计算比直接构造LoG核便宜，SIFT正是利用了这一点。

零交叉路线也有明显的缺陷。它不区分强边缘和弱边缘，平坦区域里的微小起伏同样产生零交叉，若不加幅值阈值，结果会出现大量闭合的细碎曲线；拉普拉斯算子各向同性，丢失了边缘方向；在两条边缘相交的角点附近，零交叉曲线会偏离真实位置。这些问题说明，仅凭零交叉还不足以给出干净而准确的边缘。

\subsection{Canny检测器}
\label{sec:vision-features-canny}

前两条路线各有长处：一阶路线保留了强度与方向，二阶路线给出不依赖阈值的位置。但二者都是先选定一个算子，再看它表现如何。Canny的做法相反\scite{vision-canny1986}：先写出“好的边缘检测器”应当满足的准则，再在这些准则下求最优算子。这样得到的算子不仅有理论依据，还明确揭示了噪声抑制与定位精度之间无法同时兼得的关系。

Canny在一维上建立模型：边缘是幅值为$A$的阶跃$A\,u(x)$，叠加了单位长度功率为$n_0^2$的高斯白噪声；检测器是定义在$[-W,W]$上的有限冲激响应滤波器$f$，在滤波输出的局部极大处报告边缘。他提出三条准则。

第一条是检测准则：真实边缘应以高概率被检出，噪声不应产生虚假边缘。两类错误概率都随输出信噪比单调下降，所以只需最大化信噪比。在边缘中心$x=0$处，滤波器对阶跃的响应为$\int_{-W}^{W}A\,u(-x)f(x)\,\mathrm dx=A\int_{-W}^{0}f(x)\,\mathrm dx$；白噪声通过滤波器后的标准差为$n_0\bigl(\int_{-W}^{W}f^2\bigr)^{1/2}$，原因与上一节离散情形相同，只是求和换成了积分。于是
\begin{equation}
  \Sigma(f)=\frac{A\,\bigl\lvert\int_{-W}^{0}f(x)\,\mathrm dx\bigr\rvert}{n_0\sqrt{\int_{-W}^{W}f^2(x)\,\mathrm dx}}.
  \label{eq:vision-features-canny-snr}
\end{equation}

第二条是定位准则：报告的边缘位置应尽量靠近真实位置。设输出$H(x)=H_G(x)+H_n(x)$由信号响应和噪声响应组成，报告位置$x_0$是$H$的局部极大点，满足$H_G'(x_0)+H_n'(x_0)=0$。信号响应为$H_G(x)=A\int_{-W}^{W}u(x-t)f(t)\,\mathrm dt=A\int_{-W}^{x}f(t)\,\mathrm dt$，所以$H_G'(x)=Af(x)$、$H_G''(x)=Af'(x)$；对奇函数形状的检测算子$f(0)=0$，真实边缘处$H_G'(0)=0$。在$0$附近一阶泰勒展开得$H_G'(x_0)\approx H_G''(0)x_0=Af'(0)x_0$，于是$x_0\approx-H_n'(x_0)/(Af'(0))$；噪声导数$H_n'$相当于噪声通过滤波器$f'$，其方差为$n_0^2\int f'^2$。因此定位误差的均方值约为$n_0^2\int f'^2/(A^2f'(0)^2)$，取其倒数的平方根作为定位指标：
\begin{equation}
  \Lambda(f)=\frac{A\,\lvert f'(0)\rvert}{n_0\sqrt{\int_{-W}^{W}f'^2(x)\,\mathrm dx}}.
  \label{eq:vision-features-canny-localization}
\end{equation}

这两个指标存在冲突。把滤波器在空间上拉宽$a$倍，即$f_a(x)=f(x/a)$，换元可得$\int_{-aW}^{0}f_a=a\int_{-W}^0 f$、$\int f_a^2=a\int f^2$，所以$\Sigma(f_a)=\sqrt a\,\Sigma(f)$；又$f_a'(x)=f'(x/a)/a$，所以$f_a'(0)=f'(0)/a$、$\int f_a'^2=\int f'^2/a$，于是$\Lambda(f_a)=\Lambda(f)/\sqrt a$。加宽滤波器提高信噪比、降低定位精度，二者的乘积$\Sigma\Lambda$与宽度无关，只取决于滤波器的形状。Canny因此以$\Sigma\Lambda$作为形状的优化目标，宽度则留给使用者按噪声水平选择。

第三条是单一响应准则：一个边缘只应产生一个响应。仅用前两条准则，最优解会在边缘附近产生多个靠得很近的噪声极大值。Canny利用平稳高斯过程的一个结论：滤波器对白噪声的输出中，相邻两个局部极大之间的平均距离为
\[
  x_{\max}=2\pi\Bigl(\frac{\int f'^2(x)\,\mathrm dx}{\int f''^2(x)\,\mathrm dx}\Bigr)^{1/2}.
\]
这个结论来自Rice关于高斯过程零交叉平均密度的公式，应用于$H_n'$的零交叉即为$H_n$极值的密度；它要求噪声平稳且服从高斯分布。Canny规定$x_{\max}=kW$，用常数$k$控制响应的稀疏程度。在这一约束下最大化$\Sigma\Lambda$，需要用变分法求解，Canny得到的最优阶跃边缘算子由指数与三角函数组合而成，形状与高斯一阶导数非常接近；用后者代替只损失少量性能，却能利用高斯核的可分离性高效计算。

可以直接验证高斯导数的指标。取$f(x)=-x\exp(-x^2/(2\sigma^2))$，并令$W\to\infty$。记$\int x^{2m}\mathrm e^{-x^2/\sigma^2}\mathrm dx$在$m=0,1,2$时分别为$\sqrt\pi\sigma$、$\sqrt\pi\sigma^3/2$、$3\sqrt\pi\sigma^5/4$（高斯积分的标准结果）。于是$\int_{-\infty}^0f=\sigma^2$，$\int f^2=\sqrt\pi\sigma^3/2$，$f'(0)=-1$，$\int f'^2=\int(1-x^2/\sigma^2)^2\mathrm e^{-x^2/\sigma^2}\mathrm dx=\tfrac34\sqrt\pi\sigma$。代入两式，
\[
  \Sigma=\frac{A}{n_0}\Bigl(\frac{2}{\sqrt\pi}\Bigr)^{1/2}\sigma^{1/2},\qquad
  \Lambda=\frac{A}{n_0}\Bigl(\frac{4}{3\sqrt\pi}\Bigr)^{1/2}\sigma^{-1/2},
\]
因此
\[
  \Sigma\Lambda=\frac{A^2}{n_0^2}\sqrt{\frac{8}{3\pi}}\approx0.92\,\frac{A^2}{n_0^2}.
\]
$\sigma$的作用正如前面的缩放分析所预期：它在$\Sigma$与$\Lambda$之间交换，却不改变乘积。

把一维结论推广到二维，就得到通常所说的Canny检测器。它的输出不只是一个滤波结果，还包括两步后处理，分别回答一阶路线留下的“边缘过宽”和“阈值两难”两个问题。

\begin{algorithm}[Canny边缘检测]
  \label{alg:vision-features-canny}
  \begin{algorithmic}[1]
    \Require 灰度图像$I$，平滑尺度$\sigma$，高、低阈值$T_{\mathrm{high}}>T_{\mathrm{low}}$
    \Ensure 单像素宽的二值边缘图$E$
    \State 计算$L=G_{\sigma}*I$及其梯度$L_x,L_y$（等价于与高斯导数核卷积）
    \State 按式\eqref{eq:vision-features-gradient-magnitude}计算梯度幅值$M$与方向$\alpha$
    \For{每个像素$p$}
      \State 沿$\alpha(p)$及其反方向取两个相邻位置的幅值$M_+,M_-$（落在像素之间时线性插值）
      \If{$M(p)<M_+$或$M(p)<M_-$}
        \State 令$M(p)\gets0$ \Comment{非极大值抑制}
      \EndIf
    \EndFor
    \State 标记$M\geq T_{\mathrm{high}}$的像素为强边缘，$T_{\mathrm{low}}\leq M<T_{\mathrm{high}}$的像素为弱边缘
    \State 从每个强边缘出发沿8邻接遍历，把可达的弱边缘改为强边缘 \Comment{滞后阈值}
    \State 令$E$为全部强边缘，其余像素置零
  \end{algorithmic}
\end{algorithm}

\term{非极大值抑制}（non-maximum suppression, NMS）在这里沿梯度方向比较：边缘的真实位置是剖面上梯度最大的点，它两侧的像素虽然也有较大梯度，但不是极大，应当去掉。这样每条边缘在垂直于自身的方向上只剩一个像素宽。比较只沿梯度方向进行而不沿边缘方向，因为沿边缘方向相邻的像素都属于同一条边缘，它们的幅值本应相近。

\term{滞后阈值}（hysteresis thresholding）用两个阈值区分三种像素。高于高阈值的像素足够可靠，直接作为边缘；低于低阈值的直接舍弃；介于二者之间的弱边缘，只有当它通过一串弱边缘与某个强边缘相连时才保留。它依据的假设是：真实轮廓是连续的，轮廓上即使有一段对比度降低，也会与同一轮廓上对比度较高的部分相连；而噪声点通常孤立出现。实际中高阈值常取低阈值的两到三倍。

图\ref{fig:vision-features-canny-steps}给出两步后处理的教学算例。左图中梯度方向为水平，中心幅值$120$大于左右两侧的$80$与$100$，因此保留；若中心值为$90$而右侧为$120$，中心就会被抑制，边缘位置让给右侧像素。右图是非极大值抑制之后沿一条轮廓排列的八个候选点，幅值依次为$15,42,38,25,18,33,60,12$，取$T_{\mathrm{high}}=40$、$T_{\mathrm{low}}=20$。强边缘是位置$1$和$6$；位置$2,3$的$38,25$通过位置$1$与强边缘相连而保留；位置$4$的$18$低于低阈值，在那里轮廓断开；位置$5$的$33$与强边缘$60$相邻而保留；位置$0,7$被舍弃。若改用单一阈值$40$，只剩两个孤立点；改用单一阈值$20$，则会把同样处在$20$到$40$之间、却与任何强边缘都不相连的噪声点也保留下来。滞后阈值用连通关系把“幅值”和“上下文”结合起来。

\begin{figure}[htbp]
  \centering
  % Canny 两步后处理的教学例；数值与正文算例一致。左：梯度方向为水平时的非极大值抑制；
% 右：沿一条轮廓的八个候选点与双阈值（T_high=40，T_low=20）。宽约 98mm。
\begin{tikzpicture}[x=1mm,y=1mm,font=\figurefont\small,
  cell/.style={draw=FigureStroke,line width=.5pt,minimum size=8mm,inner sep=0pt},
  cmp/.style={cell,fill=FigureBlueFill,draw=FigureBlue,line width=.8pt},
  ctr/.style={cell,fill=FigureYellowFill,draw=FigureYellow,line width=.9pt},
  lab/.style={font=\figurefont\footnotesize,text=FigureMuted}]
  % (a) NMS
  \foreach \r/\row in {0/{30,45,28},1/{80,120,100},2/{35,50,40}}{
    \foreach \v [count=\c from 0] in \row{
      \node[cell] at (\c*8,-\r*8) {\v};
    }
  }
  \node[cmp] at (0,-8) {80};
  \node[ctr] at (8,-8) {120};
  \node[cmp] at (16,-8) {100};
  \draw[-{Stealth[length=2mm]},FigureStroke,line width=.9pt] (-2,-21)--(18,-21);
  \node[lab,anchor=north] at (8,-22.5) {梯度方向};
  \node at (8,8) {(a) 非极大值抑制};
  % (b) hysteresis
  \begin{scope}[xshift=33mm,yshift=-26mm,y=.42mm]
    \foreach \v [count=\k from 0] in {15,42,38,25,18,33,60,12}{
      \pgfmathsetmacro{\keep}{(\k==1||\k==2||\k==3||\k==5||\k==6)?1:0}
      \ifdim\keep pt>0pt
        \fill[FigureBlueFill,draw=FigureBlue,line width=.7pt] (\k*7.4-3,0) rectangle (\k*7.4+3,\v);
      \else
        \draw[FigureMuted,line width=.6pt] (\k*7.4-3,0) rectangle (\k*7.4+3,\v);
      \fi
      \node[font=\figurefont\footnotesize,anchor=south] at (\k*7.4,\v) {\v};
      \node[font=\figurefont\footnotesize,anchor=north] at (\k*7.4,-1) {$\k$};
    }
    \draw[FigureStroke,line width=.6pt] (-5,0)--(57,0);
    \draw[FigureRed,dashed,line width=.7pt] (-5,40)--(57,40) node[right,font=\figurefont\footnotesize,text=FigureStroke] {$T_{\mathrm{high}}$};
    \draw[FigureYellow,dashed,line width=.7pt] (-5,20)--(57,20) node[right,font=\figurefont\footnotesize,text=FigureStroke] {$T_{\mathrm{low}}$};
    \node at (26,82) {(b) 滞后阈值};
  \end{scope}
\end{tikzpicture}

  \caption{Canny检测器两步后处理的教学例。(a) 梯度方向为水平时，中心像素（黄框）只与梯度方向上的两个邻点（蓝框）比较，$120$为极大而保留。(b) 沿一条轮廓排列的八个候选点，红色虚线为高阈值、黄色虚线为低阈值；填色柱为最终保留的边缘，空心柱被舍弃。位置$4$低于低阈值，使轮廓在该处断开。数值为教学构造。}
  \label{fig:vision-features-canny-steps}
\end{figure}

Canny检测器在噪声适中、边缘近似阶跃的图像上表现稳定，至今仍是许多系统中的默认选择。它的局限同样来自模型假设。阶跃模型不描述细线与屋脊形边缘，车道线这样的细亮条会产生两条平行的边缘，而不是一条中心线；在多条边缘交汇的角点和T形结，高斯平滑使梯度方向混乱，非极大值抑制往往在交汇处留下缺口；$\sigma$和两个阈值都要人工设置，对整幅图使用同一组参数，难以同时适应近处的粗边缘和远处的细边缘；它只看灰度变化，无法区分物体轮廓与路面纹理、阴影边界。最后一个问题是根本性的：哪些灰度变化算“有意义的边界”，取决于场景语义，而这无法由局部的平滑和阈值决定。学习式方法用人工标注的物体轮廓作为监督，让网络学习哪些变化应当报告为边缘，\ref{chap:vision-image-recognition}章中的分割网络输出的类别边界就是这种语义边界。卷积网络的第一层通常也学到类似高斯导数的方向敏感滤波器，Canny的平滑求导结构在这里以可学习权重的形式保留了下来。

\subsection{Hough变换}
\label{sec:vision-features-hough}

边缘检测给出的是一组零散的边缘点，要得到车道线的方程，还需要回答哪些点共线。车道线上的边缘点可能被车辆遮挡而断成几段，又夹杂着路面裂缝和其他车辆的边缘点。逐对检验所有点对、再统计每条候选直线上的点数，计算量随点数平方增长，并且仍要处理断裂与干扰。\term{Hough变换}（Hough transform）换了一个角度：不从点出发找直线，而是让每个点为所有可能经过它的直线投票，得票最多的直线就是被最多点支持的直线。

直线的斜截式$y=ax+b$在竖直线处斜率无界，参数空间不能有限离散化。Duda与Hart改用法线式\scite{vision-duda1972}
\begin{equation}
  \rho=x\cos\theta+y\sin\theta,\qquad \theta\in[0,\pi),\ \rho\in[-\rho_{\max},\rho_{\max}],
  \label{eq:vision-features-hough-line}
\end{equation}
其中$\theta$是直线法向量与$x$轴的夹角，$\rho$是原点到直线的有向距离，$\rho_{\max}$取图像对角线长。每条直线对应唯一的$(\rho,\theta)$，两个参数都有界。固定一个图像点$(x_i,y_i)$，满足式\eqref{eq:vision-features-hough-line}的所有$(\rho,\theta)$构成参数平面上的一条正弦曲线，它代表经过该点的全部直线。若若干点共线于参数为$(\rho_0,\theta_0)$的直线，那么它们的正弦曲线都经过点$(\rho_0,\theta_0)$。“找共线点”就转化为“找许多曲线的公共交点”。

\begin{figure}[htbp]
  \centering
  % Hough 变换教学例：三点 (1,3),(2,2),(3,1) 位于直线 x+y=4 上；右图为
% rho=x cos(theta)+y sin(theta) 的三条正弦曲线，交于 (45 度, 2.83)。宽约 104mm。
\begin{tikzpicture}[x=1mm,y=1mm,font=\figurefont\small,
  ax/.style={-{Stealth[length=1.8mm]},FigureStroke,line width=.6pt},
  lab/.style={font=\figurefont\footnotesize}]
  % image plane, 1 unit = 8mm
  \begin{scope}[x=8mm,y=8mm]
    \draw[ax] (0,0)--(4.4,0) node[right] {$x$};
    \draw[ax] (0,0)--(0,4.4) node[above] {$y$};
    \foreach \t in {1,2,3,4}{\draw[FigureStroke] (\t,0)--(\t,-.08) node[below,lab] {\t};
      \draw[FigureStroke] (0,\t)--(-.08,\t) node[left,lab] {\t};}
    \draw[FigureBlue,line width=.9pt] (0.3,3.7)--(3.7,0.3);
    \draw[-{Stealth[length=1.8mm]},FigureYellow,line width=.9pt] (0,0)--(2,2);
    \node[lab,anchor=north west] at (1.05,1.0) {$\rho$};
    \draw[FigureYellow,line width=.6pt] (0.7,0) arc[start angle=0,end angle=45,radius=0.7];
    \node[lab] at (0.95,0.32) {$\theta$};
    \foreach \x/\y in {1/3,2/2,3/1}{\fill[FigureRed] (\x,\y) circle (0.9mm);}
    \node at (2.2,4.9) {图像平面};
  \end{scope}
  % parameter plane: theta 0..180 deg -> 0..60mm, rho -3.5..4.5 -> scale 6mm
  \begin{scope}[xshift=44mm,yshift=22mm]
    % theta axis at rho=-3.5 (bottom), rho=0 as grid line
    \draw[FigureGrid,line width=.5pt] (0,0)--(50,0);
    \draw[ax] (0,-21)--(54,-21) node[right] {$\theta$};
    \draw[ax] (0,-21)--(0,24) node[above] {$\rho$};
    \foreach \d in {0,45,90,135,180}{\draw[FigureStroke] ({\d/3.6},-21)--({\d/3.6},-22) node[below,lab] {$\d^\circ$};}
    \foreach \r in {-3,0,3}{\draw[FigureStroke] (0,{6*\r})--(-1,{6*\r}) node[left,lab] {$\r$};}
    \tikzset{houghcurve1/.style={FigureBlue},houghcurve2/.style={FigureStroke,dashed},houghcurve3/.style={FigureRed}}
    \foreach \px/\py/\c in {1/3/houghcurve1,2/2/houghcurve2,3/1/houghcurve3}{
      \draw[\c,line width=.9pt,domain=0:180,samples=60,smooth]
        plot ({\x/3.6},{6*(\px*cos(\x)+\py*sin(\x))});
    }
    \node[lab,anchor=west,text=FigureBlue] at (50.5,-6) {$(1,3)$};
    \node[lab,anchor=west,text=FigureStroke] at (50.5,-12) {$(2,2)$};
    \node[lab,anchor=west,text=FigureRed] at (50.5,-18) {$(3,1)$};
    \fill[FigureStroke] (12.5,{6*2.8284}) circle (1mm);
    \draw[FigureMuted,dashed,line width=.5pt] (12.5,-21)--(12.5,{6*2.8284});
    \node[lab,anchor=south] at (12.5,{6*2.8284+1.5}) {$(45^\circ,\,2.83)$};
    \node at (25,30) {参数平面};
  \end{scope}
\end{tikzpicture}

  \caption{Hough变换的点—线对偶。左：三个点$(1,3),(2,2),(3,1)$位于直线$x+y=4$上，该直线的法向角$\theta=45^\circ$，原点到直线的距离$\rho=4/\sqrt2\approx2.83$。右：每个点按式\eqref{eq:vision-features-hough-line}在参数平面上画出一条正弦曲线，曲线右端标出对应的点；三条曲线交于$(45^\circ,2.83)$。图中坐标轴按教学需要选取，$y$轴取向上。}
  \label{fig:vision-features-hough}
\end{figure}

计算时把参数平面离散为二维累加器数组，$\theta$取$n_{\theta}$个值，$\rho$按步长$\Delta\rho$量化。每个边缘点对每个$\theta$算出$\rho$，在对应格子上加一票；全部点投完后，找累加器中的局部极大值，并用非极大值抑制去掉相邻格子中的重复峰。以图\ref{fig:vision-features-hough}的三个点为例，取$\theta\in\{0^\circ,45^\circ,90^\circ,135^\circ\}$、$\Delta\rho=1$并就近取整：点$(1,3)$给出$\rho=1,2.83,3,1.41$，点$(2,2)$给出$2,2.83,2,0$，点$(3,1)$给出$3,2.83,1,-1.41$。只有$\theta=45^\circ$、$\rho\approx3$的格子得到三票，其余格子都只得到一票（例如$\theta=0^\circ$时三个点分别落在$\rho=1,2,3$三个格子），峰值给出直线$x\cos45^\circ+y\sin45^\circ\approx2.83$。

投票的计算量是边缘点数乘以$n_{\theta}$，与点对数无关。Hough变换对断裂天然不敏感：一条直线上的点无论是否连续，都投给同一个格子。它对少量干扰点也较稳健，因为干扰点的票分散在许多格子中，不会集中到某一格。实践中有两个常用改进。一是用梯度方向限制投票：边缘点的梯度方向就是所在直线的法向，因此只需在$\alpha$附近的少数$\theta$上投票，既降低计算量又减少错误票。二是按梯度幅值加权投票，使强边缘贡献更多。

同样的方法可推广到任何可参数化的曲线。圆$(x-a)^2+(y-b)^2=r^2$有三个参数，累加器变成三维；若已知边缘点的梯度方向，圆心必在过该点、沿梯度方向的直线上，对每个半径$r$只需在两个候选圆心处投票，计算量大大减少。参数越多，累加器的维数越高，存储与计算按指数增长，所以Hough变换主要用于直线、圆这类低维形状。

累加器的分辨率决定了精度与稳健性的取舍。格子太粗，相近的两条直线落进同一格，定位粗糙；格子太细，同一条略有弯曲或带噪声的直线的票分散到相邻格子，峰值不再突出。透视使车道线在远处略有弯曲，在细格子下峰值会被拉成一条带。Hough变换也不能告诉我们直线在哪里开始、在哪里结束，需要回到图像中沿峰值直线检查边缘点，才能得到线段的端点。它与\ref{sec:vision-recognition-ransac}节的RANSAC都用于在大量干扰中找出被多数点支持的模型：Hough变换穷举离散化的参数空间，适合参数维数低的情形；RANSAC随机抽取最小样本生成候选模型，适合参数维数较高、但单个模型可由少数点确定的情形。投票汇总局部证据的思想在SIFT的物体识别流程中也被采用\scite{vision-lowe2004}。

表\ref{tab:vision-features-edge-compare}在同一街景输入下比较本节的四种方法。

\begin{table*}[tbp]
  \centering
  \small
  \begin{tabularx}{\linewidth}{@{}lXXXX@{}}
    \toprule
    比较项 & 一阶梯度阈值 & LoG零交叉 & Canny & Hough变换 \\
    \midrule
    输入 & 灰度图 & 灰度图 & 灰度图 & 边缘点（可带方向） \\
    输出 & 带强度与方向的粗边缘 & 闭合的零交叉曲线 & 单像素宽、连接的边缘 & 直线或圆的参数 \\
    关键设计 & 差分核与平滑核的组合 & 先高斯平滑再求拉普拉斯 & 三准则导出的高斯导数、非极大值抑制、滞后阈值 & 点—参数对偶与投票 \\
    需设参数 & 阈值 & 尺度$\sigma$、零交叉强度阈值 & $\sigma$、两个阈值 & 参数离散步长、峰值阈值 \\
    典型失效 & 边缘过宽、阈值两难 & 碎曲线多、无方向、角点偏移 & 交汇处缺口、细线成双边、不分语义 & 参数维数高时不可行、峰值分散 \\
    \bottomrule
  \end{tabularx}
  \caption{本节四种方法的比较。前三种回答哪些像素是边缘，第四种把边缘点组织成几何对象，通常接在Canny之后使用。}
  \label{tab:vision-features-edge-compare}
\end{table*}

\section{局部特征的检测与描述}
\label{sec:vision-features-local-features}

边缘告诉我们灰度在哪里变化，但边缘上的点彼此相似，沿着边缘滑动看不出区别。现在设想要在两幅街景照片之间找到同一块招牌：第二幅照片在另一位置拍摄，招牌更小、略有旋转，光照也不同。要建立两幅图之间的对应，需要满足两个要求。第一是可重复检测：同一个物理点在两幅图中都应被选中，并且选中的位置、尺度和方向要随图像变换一起变化。第二是可区分描述：每个被选中的点要附上一个向量，同一物理点的两个向量彼此接近，不同点的向量相距较远，并且这个向量不受视角、尺度与光照变化的显著影响。满足这两点的点及其向量称为\term{局部特征}（local feature）；前者由检测子给出，后者由\term{描述子}（descriptor）给出。

局部特征不只用于匹配。识别行人时，我们关心的是窗口内梯度方向的整体分布，而不是某几个关键点；区分路面和草地时，我们关心的是纹理的统计规律。这三种需求对应本节的不同路线，它们的差别在于在哪里取特征、统计什么。Harris角点只解决检测：找出在两个方向上都有灰度变化、因而能精确定位的点，但它不处理尺度变化。SIFT同时解决检测与描述，并把尺度纳入检测：在尺度空间中寻找极值，按局部梯度确定主方向，再在归一化的邻域中统计梯度方向直方图。HOG放弃关键点，在整个检测窗口上稠密地统计梯度方向，用于描述一类对象的整体形状。LBP与灰度共生矩阵也是稠密统计，但统计对象是像素之间的灰度大小关系或共现频率，用于描述纹理。

\subsection{Harris角点}
\label{sec:vision-features-harris}

先问什么样的点适合定位。把一个小窗口放在图像上并稍作移动，比较移动前后窗口内容的差别。在平坦区域，向任何方向移动内容都几乎不变；在边缘上，沿边缘方向移动内容不变，只有垂直于边缘移动才变化；在角点处，向任何方向移动内容都明显变化。只有最后一种点能在两个方向上都被唯一确定，这样的点称为\term{角点}（corner）。Harris与Stephens把这一观察写成二次型\scite{vision-harris1988}。

设窗口权重为$w(x,y)\geq0$（常取高斯函数），窗口平移$(u,v)$后的加权平方差为
\[
  E(u,v)=\sum_{x,y}w(x,y)\bigl[I(x+u,y+v)-I(x,y)\bigr]^2.
\]
对小位移作一阶泰勒展开$I(x+u,y+v)\approx I(x,y)+I_xu+I_yv$，代入得
\begin{equation}
  E(u,v)\approx\begin{pmatrix}u&v\end{pmatrix}\bm M\begin{pmatrix}u\\v\end{pmatrix},\qquad
  \bm M=\sum_{x,y}w(x,y)\begin{pmatrix}I_x^2&I_xI_y\\I_xI_y&I_y^2\end{pmatrix}.
  \label{eq:vision-features-structure-tensor}
\end{equation}
矩阵$\bm M$称为\term{结构张量}（structure tensor），也叫二阶矩矩阵。它是各像素梯度外积$\nabla I\,\nabla I^{\top}$的加权和；每个外积都是对称半正定矩阵，非负加权和保持这一性质，所以$\bm M$对称半正定，有两个非负特征值$\lambda_1\geq\lambda_2\geq0$及正交的特征向量。

特征值的含义来自对称矩阵的瑞利商：对单位向量$\bm e$，$\bm e^{\top}\bm M\bm e$介于$\lambda_2$与$\lambda_1$之间，在对应特征向量处取到两端，可参考\BookChapterRef{chap:linear-analysis}{卷一“特征值与谱理论”章}。因此$\lambda_1$是窗口沿变化最快方向移动单位距离时的平方差，$\lambda_2$是沿变化最慢方向的平方差。平坦区域$\lambda_1\approx\lambda_2\approx0$；边缘处$\lambda_1$大而$\lambda_2\approx0$，最慢方向就是边缘方向；角点处两个特征值都大。判断角点只需要检查$\lambda_2$是否足够大。

Harris没有直接求特征值，而是用迹与行列式构造响应函数
\begin{equation}
  R=\det\bm M-\kappa\,(\operatorname{tr}\bm M)^2=\lambda_1\lambda_2-\kappa(\lambda_1+\lambda_2)^2,
  \label{eq:vision-features-harris}
\end{equation}
其中$\kappa$是经验常数，实践中常取$0.04$至$0.06$。对$2\times2$矩阵，行列式和迹只需几次乘加。$R$在角点处为大的正值，在边缘处为负值（因为$\lambda_2\approx0$时$R\approx-\kappa\lambda_1^2$），在平坦区域接近零。Shi与Tomasi后来直接以$\min(\lambda_1,\lambda_2)$作为判据，用于选择适合跟踪的点\scite{vision-shi1994good}；两种判据在大多数情形下给出相近的结果。

用一个算例比较三种情形。设窗口内有$8$个像素，$w\equiv1$，梯度幅值都是$10$。在平坦区域，所有梯度为零，$\bm M=\bm0$，$R=0$。在一条竖直边缘上，8个像素的梯度都是$(10,0)$，$\bm M=\operatorname{diag}(800,0)$，于是$\det\bm M=0$、$\operatorname{tr}\bm M=800$，取$\kappa=0.04$得$R=-0.04\times800^2=-25600$。在一个直角角点处，4个像素位于竖直边上、梯度为$(10,0)$，另4个位于水平边上、梯度为$(0,10)$，$\bm M=\operatorname{diag}(400,400)$，$R=160000-0.04\times800^2=134400$。三种情形的$R$分别为零、负值和大的正值。注意边缘和角点的$\operatorname{tr}\bm M$相同，只看总梯度能量无法区分二者；区分它们的是能量在两个方向上的分配。

\begin{algorithm}[Harris角点检测]
  \label{alg:vision-features-harris}
  \begin{algorithmic}[1]
    \Require 灰度图像$I$，求导尺度$\sigma_D$，积分尺度$\sigma_I$，常数$\kappa$，阈值$T$
    \Ensure 角点位置集合
    \State 用尺度为$\sigma_D$的高斯导数核计算$I_x,I_y$
    \State 逐像素计算$I_x^2,\ I_y^2,\ I_xI_y$
    \State 用尺度为$\sigma_I$的高斯核分别平滑这三幅图，得到每个像素处$\bm M$的三个元素
    \State 按式\eqref{eq:vision-features-harris}计算响应图$R$
    \State 保留$R>T$且在$3\times3$邻域内为极大的像素
  \end{algorithmic}
\end{algorithm}

Harris角点有哪些不变性？若图像绕某点旋转，梯度随之旋转，$\bm M$变为$\bm Q\bm M\bm Q^{\top}$（$\bm Q$为旋转矩阵），特征值不变，所以检测结果对旋转不变。灰度加上常数不改变梯度，因而不影响$\bm M$。灰度乘以常数$c$使$\bm M$乘以$c^2$、$R$乘以$c^4$，固定阈值$T$下检出的点数会随对比度变化，这是实际使用中需要按图像调整阈值的原因。Harris角点对尺度变化不具备不变性：同一个圆角，在高分辨率图像中，窗口只覆盖弧线的一小段，看起来像边缘；在低分辨率图像中，整个圆角落入窗口，才呈现为角点。招牌在两幅照片中大小不同，同一个窗口尺寸不可能在两幅图中都选中同一个点。这就引出下一小节的尺度空间。

结构张量在后面的章节中还会出现。\ref{sec:vision-videofound-lucas-kanade}节的Lucas–Kanade光流用同一个矩阵求解局部运动：矩阵接近奇异时只能确定垂直于边缘的运动分量，这正是孔径问题；Harris角点在那里也就是可以可靠跟踪的点。\ref{sec:vision-recognition-superpoint}节的SuperPoint则用网络学习兴趣点检测，并以合成角点作为初始监督。

\subsection{SIFT}
\label{sec:vision-features-sift}

\term{尺度不变特征变换}（scale-invariant feature transform, SIFT）要解决的问题是：在招牌大小不同、方向不同、光照不同的两幅照片中，独立地检测出同一组点，并为每个点给出可以直接比较的描述\scite{vision-lowe2004}。它的思路是把尺度当作和位置一样的待检测量：不预先规定窗口大小，而是在一组连续尺度上观察图像，在位置与尺度构成的三维空间中寻找响应的极值；每个极值同时给出位置和该点的“自然尺度”。随后按局部梯度确定主方向，在按尺度缩放、按方向旋转后的邻域中计算描述子，于是描述子对尺度和旋转不再敏感。

先说明怎样在多个尺度上观察一幅图像。\term{尺度空间}（scale space）把图像$I$嵌入一族逐渐平滑的图像
\begin{equation}
  L(x,y,\sigma)=G_{\sigma}*I,\qquad \sigma>0.
  \label{eq:vision-features-scale-space}
\end{equation}
高斯核被选作尺度空间的核，有两条理由。其一是半群性质：$G_{\sigma_1}*G_{\sigma_2}=G_{\sqrt{\sigma_1^2+\sigma_2^2}}$，两次平滑等于一次更大尺度的平滑，所以较粗尺度的图像可以由较细尺度的图像继续平滑得到，不必每次从原图重算。这一性质来自两个独立高斯随机向量之和仍服从高斯分布、方差相加。其二，$L$作为$t=\sigma^2/2$的函数满足热传导方程$\partial L/\partial t=\nabla^2L$，即上一节已经验证的$\partial G_{\sigma}/\partial\sigma=\sigma\nabla^2G_{\sigma}$。在一维情形下，热方程的解随时间推移不会产生新的局部极值，平滑只会合并或消除结构，而不会凭空制造结构。尺度空间理论表明，在线性、平移与旋转不变并满足这种“不产生新结构”的公理下，高斯核是唯一的选择，Lindeberg对这些结果有系统的整理\scite{vision-lindeberg1998}；二维中局部极值的个数并不严格随尺度单调减少，“不产生新结构”在二维要以更弱的形式表述。

尺度空间解决了“在每个尺度上看”，还需要回答“哪个尺度才是这个结构的尺度”。答案是在尺度方向上找导数响应的极值，但响应必须经过归一化。直接看LoG响应$\nabla^2L$，平滑越强，图像越平坦，响应随$\sigma$单调下降，极值永远在最小尺度处。Lindeberg提出乘以$\sigma^2$得到尺度归一化的LoG，$\sigma^2\nabla^2L$，补偿平滑造成的幅值衰减\scite{vision-lindeberg1998}。

一个斑点的算例说明归一化后的极值恰好落在斑点自身的尺度上。设图像是中心在原点、宽度为$s_0$的高斯斑点$I(x,y)=\exp(-(x^2+y^2)/(2s_0^2))=2\pi s_0^2\,G_{s_0}(x,y)$。由半群性质，$L=2\pi s_0^2\,G_{\tau}$，其中$\tau^2=s_0^2+\sigma^2$。由式\eqref{eq:vision-features-log}，$\nabla^2G_{\tau}$在原点的值为$-2\tau^2/\tau^4\cdot1/(2\pi\tau^2)=-1/(\pi\tau^4)$。于是中心处的归一化响应为
\[
  \sigma^2\nabla^2L(0,0,\sigma)=-\frac{2s_0^2\sigma^2}{(s_0^2+\sigma^2)^2}.
\]
令$t=\sigma^2$，对$t/(s_0^2+t)^2$求导得$(s_0^2-t)/(s_0^2+t)^3$，在$t=s_0^2$处由正变负，所以响应的绝对值在$\sigma=s_0$处最大，最大值为$1/2$。取$s_0=4$，$\sigma=2,4,8$时绝对值依次为$0.32,0.50,0.32$。若斑点在另一幅图中放大一倍，$s_0=8$，极值随之移到$\sigma=8$。检测出的尺度与图像一起缩放，这就是尺度不变检测的含义。

SIFT用上一节的DoG近似尺度归一化的LoG。把上一节得到的近似式两边与$I$卷积，得
\begin{equation}
  D(x,y,\sigma)=L(x,y,k\sigma)-L(x,y,\sigma)\approx(k-1)\,\sigma^2\nabla^2L.
  \label{eq:vision-features-dog}
\end{equation}
系数$k-1$对所有尺度相同，不影响极值位置，因此DoG本身就带有尺度归一化。计算上，DoG只是相邻两幅高斯平滑图像相减，而高斯平滑图像本来就需要计算。

尺度空间按组（octave）组织，每组覆盖尺度翻倍的范围，分成$s$个间隔，相邻尺度之比$k=2^{1/s}$。为了在每组的$s$个尺度上都能与上下两层比较，每组要生成$s+3$幅高斯图像、相减得到$s+2$幅DoG图像。下一组由本组尺度为$2\sigma$的那幅图像隔点下采样得到，这与\ref{sec:vision-transforms-pyramid}节的高斯金字塔相同，下采样之前的平滑已足以抑制混叠。Lowe通过重复检测率实验选择$s=3$，每组基准尺度$\sigma=1.6$；并假设原图已有$0.5$的模糊，先用线性插值把图像放大一倍再构造第一组，这使稳定关键点数量增加近四倍\scite{vision-lowe2004}。图\ref{fig:vision-features-sift-scale-space}画出了两组的结构。

\begin{figure*}[tbp]
  \centering
  % SIFT 尺度空间构造示意：为减少层数取每组 s=2 个间隔（原文实验取 s=3），
% 每组 s+3=5 幅高斯图像、s+2=4 幅高斯差分图像；第二组边长减半。宽约 112mm。
\begin{tikzpicture}[x=1mm,y=1mm,font=\figurefont\small,
  gimg/.style={draw=FigureBlue,fill=FigureBlueFill,line width=.6pt},
  dimg/.style={draw=FigureYellow,fill=FigureYellowFill,line width=.6pt},
  link/.style={-{Stealth[length=1.8mm]},draw=FigureStroke,line width=.6pt},
  lab/.style={font=\figurefont\footnotesize,text=FigureMuted}]
  % octave 1: 5 gaussians stacked upward, size 16x10 slanted
  \foreach \k in {0,...,4}{
    \draw[gimg] (0+\k*1.5,\k*6) -- ++(16,0) -- ++(5,5) -- ++(-16,0) -- cycle;
  }
  \foreach \k in {0,...,3}{
    \draw[dimg] (40+\k*1.5,3+\k*6) -- ++(16,0) -- ++(5,5) -- ++(-16,0) -- cycle;
    \draw[link] (23+\k*1.5,5+\k*6) -- (39+\k*1.5,5.5+\k*6);
  }
  \node[lab,anchor=east] at (-1,2.5) {$\sigma$};
  \node[lab,anchor=east] at (-1,26.5) {$k^4\sigma$};
  \node at (13,36) {高斯图像 $L$};
  \node at (53,33) {高斯差分 $D$};
  \node[lab] at (31,-4) {相邻两层相减};
  % octave 2 (half size)
  \begin{scope}[yshift=-24mm]
    \foreach \k in {0,...,4}{
      \draw[gimg] (0+\k*.9,\k*3.4) -- ++(8,0) -- ++(2.5,2.5) -- ++(-8,0) -- cycle;
    }
    \foreach \k in {0,...,3}{
      \draw[dimg] (40+\k*.9,1.7+\k*3.4) -- ++(8,0) -- ++(2.5,2.5) -- ++(-8,0) -- cycle;
    }
    \node[lab,anchor=north west] at (-2,-1) {第二组：由尺度为 $k^s\sigma=2\sigma$ 的一层隔点下采样得到};
  \end{scope}
  \draw[link] (3,14.5) .. controls (-7,12) and (-7,-14) .. (0.2,-16);
  % 26-neighbour
  \begin{scope}[xshift=82mm,yshift=4mm]
    \foreach \l/\yy in {0/0,1/9,2/18}{
      \foreach \i in {0,1,2}{\foreach \j in {0,1,2}{
        \draw[FigureStroke,line width=.4pt,fill=white] (\i*3.6+\l*1.2,\yy+\j*2.2) rectangle ++(3.6,2.2);
      }}
    }
    \fill[FigureRed] (3.6+1.2+1.8,9+2.2+1.1) circle (0.9mm);
    \node[lab,anchor=west] at (13,1.5) {$D(\cdot,\cdot,\sigma_{i-1})$};
    \node[lab,anchor=west] at (13,10.5) {$D(\cdot,\cdot,\sigma_{i})$};
    \node[lab,anchor=west] at (13,19.5) {$D(\cdot,\cdot,\sigma_{i+1})$};
    \node at (8,29) {26 个邻点};
  \end{scope}
\end{tikzpicture}

  \caption{SIFT的尺度空间与极值检测。每组内的高斯图像$L$按尺度从下到上排列，相邻两幅相减得到高斯差分$D$；下一组由隔点下采样得到，边长减半。右：DoG中的一个采样点（红点）与同层8个邻点及上下两层各9个点比较。为减少层数，图中取每组$s=2$个间隔，即5幅高斯图像、4幅DoG；Lowe的实验设置为$s=3$。}
  \label{fig:vision-features-sift-scale-space}
\end{figure*}

在DoG中，每个采样点与同层的8个邻点以及上下两层各9个点共26个点比较，大于全部或小于全部时成为候选关键点。离散采样上的极值并不精确：真实极值可能落在采样点之间，也可能是对比度很低、容易受噪声影响的点，或者落在边缘上、沿边缘方向定位不稳的点。SIFT对候选点做三步处理。

第一步是亚像素定位。把$D$在采样点处对$\bm x=(x,y,\sigma)^{\top}$的偏移作二阶泰勒展开，$D(\bm x)\approx D+\nabla D^{\top}\bm x+\tfrac12\bm x^{\top}\bm H\bm x$，其中梯度$\nabla D$与海森矩阵$\bm H$都用相邻采样的差分近似。令对$\bm x$的梯度为零，得极值点偏移和极值
\begin{equation}
  \hat{\bm x}=-\bm H^{-1}\nabla D,\qquad
  D(\hat{\bm x})=D+\tfrac12\nabla D^{\top}\hat{\bm x}.
  \label{eq:vision-features-sift-refine}
\end{equation}
第二式由代入$\bm H\hat{\bm x}=-\nabla D$得到：$\tfrac12\hat{\bm x}^{\top}\bm H\hat{\bm x}=-\tfrac12\hat{\bm x}^{\top}\nabla D$，与线性项合并。若$\hat{\bm x}$的某个分量超过$0.5$，说明极值更靠近另一个采样点，就移到那个点上重新展开。一维算例：三个相邻尺度上的$D$依次为$0.20,0.30,0.26$，中心差分$D'=(0.26-0.20)/2=0.03$，二阶差分$D''=0.26-2\times0.30+0.20=-0.14$，于是$\hat x=-0.03/(-0.14)\approx0.21$，极值约为$0.30+\tfrac12\times0.03\times0.21\approx0.303$。真实极值比中心采样点偏向第三个采样一侧约五分之一个间隔。

第二步剔除低对比度点：若像素值归一化到$[0,1]$，$\lvert D(\hat{\bm x})\rvert<0.03$的点被舍弃\scite{vision-lowe2004}。

第三步剔除边缘响应。DoG在边缘上也有较强的响应，但沿边缘方向的曲率很小，位置沿边缘滑动时响应几乎不变，定位不稳定。这与Harris角点区分边缘和角点的道理相同，只是这里用的是$D$在空间上的$2\times2$海森矩阵$\bm H_s$，它的两个特征值$\alpha,\beta$是两个主方向上的曲率。设$\alpha=r\beta$（$r\geq1$），则
\[
  \frac{(\operatorname{tr}\bm H_s)^2}{\det\bm H_s}=\frac{(\alpha+\beta)^2}{\alpha\beta}=\frac{(r+1)^2}{r},
\]
它只依赖两曲率之比，并随$r$增大而增大。若行列式为负，两个曲率异号，该点不是极值，直接舍弃。Lowe取$r=10$，即保留$(\operatorname{tr}\bm H_s)^2/\det\bm H_s<12.1$的点。例如$D_{xx}=-0.10$、$D_{yy}=-0.005$、$D_{xy}=0$时比值为$0.105^2/0.0005\approx22.1$，该点位于边缘上，被剔除；若$D_{yy}=-0.04$，比值为$0.14^2/0.004=4.9$，保留。整个检验只需迹和行列式，不必求特征值。

有了位置与尺度，下一步给关键点指定方向。在关键点尺度对应的高斯图像$L$上，用中心差分计算邻域内每点的梯度幅值与方向，按方向建立36个箱、每箱$10^\circ$的直方图；每个样本以梯度幅值乘以高斯窗权重计票，窗的标准差是关键点尺度的$1.5$倍。直方图的最高峰给出主方向；任何高于最高峰$80\%$的其他局部峰也各自生成一个同位置、同尺度、不同方向的关键点。峰值位置用相邻三个箱拟合抛物线插值得到。之后的描述子在以主方向为$x$轴的坐标系中计算，因而对旋转不变。

描述子在关键点周围取$16\times16$个采样点（采样间距按关键点尺度确定），梯度方向都减去主方向；以描述窗半宽为标准差的高斯函数加权，使远离中心、受配准误差影响较大的梯度贡献较小。采样点划分为$4\times4$个子区域，每个子区域统计8个方向的梯度直方图，得到$4\times4\times8=128$维向量。为避免样本跨越子区域或方向箱边界时描述子发生跳变，每个样本按到相邻箱中心的距离，在两个空间维度和一个方向维度上线性分配权重，即三线性插值。最后对光照做两步处理：先把向量归一化为单位长度，抵消对比度整体缩放（灰度加常数本就不改变梯度）；再把大于$0.2$的分量截断为$0.2$并重新归一化，减弱相机饱和或非均匀光照造成的少数大梯度的影响\scite{vision-lowe2004}。

匹配时，对一幅图中的每个描述子在另一幅图中找欧氏距离最近和次近的描述子。若只按最近距离设阈值，很难兼顾不同特征的区分度；Lowe改用最近与次近距离之比：真实匹配的最近邻明显比次近邻近，而错误匹配的最近与次近距离差不多。比值大于$0.8$时拒绝匹配，在他的实验中这一规则排除了约$90\%$的错误匹配，同时丢弃的正确匹配不到$5\%$\scite{vision-lowe2004}。例如最近与次近距离为$0.30$与$0.45$时比值约$0.67$，接受；为$0.40$与$0.44$时比值约$0.91$，拒绝，因为两个候选几乎一样像，这个匹配不可靠。窗户、栅栏等重复结构上的点常属于后一种情况。

\begin{algorithm}[SIFT特征提取]
  \label{alg:vision-features-sift}
  \begin{algorithmic}[1]
    \Require 灰度图像$I$，每组间隔数$s$，基准尺度$\sigma_0$，对比度阈值，曲率比阈值$r$
    \Ensure 关键点$(x,y,\sigma,\theta)$及其128维描述子
    \State 放大图像一倍，逐组生成$s+3$幅高斯图像与$s+2$幅DoG图像，组间隔点下采样
    \State 在DoG中找26邻域极值作为候选点
    \For{每个候选点}
      \State 按式\eqref{eq:vision-features-sift-refine}迭代求亚像素位置与尺度，偏移过大时换点
      \State 若$\lvert D(\hat{\bm x})\rvert$小于对比度阈值或曲率比检验失败，则舍弃
      \State 由36箱方向直方图确定一个或多个主方向
      \State 在按尺度缩放、按主方向旋转的邻域中计算$4\times4\times8$直方图
      \State 归一化、截断于$0.2$、再归一化
    \EndFor
  \end{algorithmic}
\end{algorithm}

SIFT的各项设计都对应一种已知的图像变化：尺度空间极值应对尺度，主方向应对旋转，梯度与归一化应对仿射光照，子区域直方图与三线性插值容忍小的位移和形变。它在平面物体和中等视角变化下的匹配效果长期是其他方法的比较基准。它的局限也来自这些设计：描述子只对相似变换（平移、旋转、缩放）严格不变，对大视角造成的透视和非平面物体的形变只能容忍一部分；弱纹理表面（白墙、天空）检测不到足够的关键点；重复纹理使比值检验拒绝大量匹配；光照造成的阴影边界会被误当作结构。计算量也较大，一幅中等大小的图像常有上千个关键点，每个都要计算128维描述子。

SIFT之后的学习方法保留了它的整体流程，只是把手工设计的部分逐一换成学习得到的部分。\ref{sec:vision-recognition-superpoint}节的SuperPoint用一个网络同时输出关键点概率图和稠密描述子，\ref{sec:vision-recognition-lightglue}节的LightGlue用注意力在两组特征之间直接推断对应关系，代替最近邻加比值检验。SIFT在多个尺度上检测、在尺度归一化的邻域中描述的想法，与卷积网络中多尺度特征金字塔的设计（\ref{sec:vision-recognition-fpn}节）也是一致的：不同大小的目标应在与其大小相称的分辨率上处理。

\subsection{HOG}
\label{sec:vision-features-hog}

现在换一个问题：判断一个$64\times128$的窗口里是否站着一个行人。行人的衣着颜色、纹理千差万别，背景也各不相同，像素值本身很难比较；但人体轮廓的形状是相对稳定的，头、肩、腿两侧的边缘方向分布大致相同。SIFT只在少数关键点上描述局部，且对旋转不变，而行人总是直立的，方向本身就是有用的信息。Dalal与Triggs提出的\term{方向梯度直方图}（histogram of oriented gradients, HOG）因此改为在整个窗口上稠密地统计梯度方向，并保留方向\scite{vision-dalal2005}。

HOG的计算分四步，Dalal与Triggs对每一步的选择都做了对比实验。

第一步计算梯度。实验表明最简单的一维中心差分$(-1,0,1)$不加预平滑效果最好；先做高斯平滑或使用更大的Sobel核反而降低检测率。对彩色图像，每个像素取三个通道中梯度幅值最大的那个通道。这与边缘检测追求抑制噪声的做法相反：行人轮廓的细节分布在细尺度上，平滑会把它们抹掉。

第二步在单元（cell）内按方向投票。窗口划分为$8\times8$像素的单元；每个单元建立9个方向箱的直方图，方向范围取$0^\circ$至$180^\circ$，即不区分梯度的正负方向。这一选择针对行人：衣服可能比背景亮也可能比背景暗，梯度的符号并不稳定。每个像素按梯度幅值投票，并在相邻两个方向箱之间线性分配。例如方向箱中心为$10^\circ,30^\circ,\ldots,170^\circ$，一个幅值为$10$、方向为$75^\circ$的像素，到中心$70^\circ$与$90^\circ$的距离分别是$5^\circ$和$15^\circ$，因此给$70^\circ$箱投$10\times15/20=7.5$票，给$90^\circ$箱投$2.5$票。投票在空间上也按位置在相邻单元之间插值，这和SIFT的三线性插值目的相同。

第三步在块（block）内做对比度归一化。光照和前景背景对比度在窗口内变化很大，阴影中的一条腿与阳光下的一条腿梯度幅值可能差几倍。把$2\times2$个相邻单元（$16\times16$像素）组成一个块，将4个单元的直方图拼成36维向量$\bm v$，用L2-Hys归一化：先$\bm v\gets\bm v/\sqrt{\lVert\bm v\rVert_2^2+\epsilon^2}$，再把大于$0.2$的分量截断为$0.2$，最后再归一化一次。块以8像素为步长滑动，相邻块重叠，每个单元出现在最多4个块中，每次按不同邻域归一化。看起来是冗余，实验却表明这种重叠提高了检测率：同一个单元在不同块中按不同的局部对比度归一化，描述子因此对局部光照变化更稳健。截断的作用可以用一个四维的教学向量看出：$\bm v=(3,4,0,0)$归一化为$(0.6,0.8,0,0)$，截断为$(0.2,0.2,0,0)$，再归一化为$(0.71,0.71,0,0)$；原先占主导的第二个分量与第一个分量拉平，单个强梯度不再支配整个块。

第四步拼接。$64\times128$的窗口有$8\times16$个单元，块在水平方向有$7$个位置、竖直方向有$15$个位置，共$105$个块，每块36维，HOG描述子为$105\times36=3780$维。检测时把这个向量交给线性支持向量机打分，具体训练与检测流程见\ref{sec:vision-features-dpm}节。

HOG的默认配置还包括：RGB色彩空间、不做伽马校正（平方根伽马压缩只带来很小的改善）、块内的高斯空间加权（标准差为8像素）。Dalal与Triggs的实验结论可以归结为三点：梯度要在细尺度上计算，方向要分得足够细（约9箱），归一化要在重叠的局部块内进行\scite{vision-dalal2005}。

HOG为什么有效，又在哪里失效？它把窗口中的形状信息转换为“每个小区域内各方向边缘的强度”，丢弃了单元内的精确位置，因此容忍小的形变和错位；归一化又抵消了大部分光照变化。但它本质上是一个刚性模板：行人抬腿、转身、被遮挡一半时，单元的直方图与训练模板的对应关系就被破坏。HOG不对旋转和尺度做任何不变处理，必须在图像金字塔上用固定大小的窗口搜索，计算量与位置数乘以尺度数成正比。

从计算结构看，HOG是一个固定的浅层卷积网络：梯度计算是与固定核的卷积，按方向分箱相当于非线性的方向选择，单元内求和是步长为8的求和池化，块归一化是局部响应归一化。\ref{sec:vision-ivm-resnet}节的ResNet保留了“卷积—非线性—汇总—归一化”的层次，但把每一步的参数都交给学习，并堆叠到几十层。

\subsection{LBP与灰度共生矩阵}
\label{sec:vision-features-texture}

前三小节描述的是边缘和形状。街景中还有另一类需要区分的信息：路面的沥青、路边的草地、建筑立面的砖墙，它们没有明确的轮廓，区别在于局部灰度起伏的重复规律，即\term{纹理}（texture）。纹理是一种统计性质，必须在一个区域上汇总才有意义。两种经典方法从不同角度描述这种统计：LBP记录每个像素与其邻域之间的大小关系，再统计这种局部模式在区域中的分布；灰度共生矩阵记录相隔固定位移的两个像素同时取各种灰度值的频率，再由这个联合分布计算对比度、能量等标量特征。

\term{局部二值模式}（local binary pattern, LBP）在中心像素$g_c$周围半径$R$的圆上等间隔取$P$个点$g_0,\ldots,g_{P-1}$（不在像素中心的点用双线性插值），把每个邻点与中心比较，得到一个$P$位二进制数\scite{vision-ojala2002}：
\begin{equation}
  \mathrm{LBP}_{P,R}=\sum_{p=0}^{P-1}s(g_p-g_c)\,2^p,\qquad
  s(z)=\begin{cases}1,&z\geq0,\\0,&z<0.\end{cases}
  \label{eq:vision-features-lbp}
\end{equation}
LBP只用到大小关系，所以对任何保持灰度次序的单调变换都不变，这包括亮度偏移、对比度缩放和伽马校正。以$P=8,R=1$为例，设$3\times3$邻域为
\[
  \begin{pmatrix}6&5&2\\7&6&1\\9&8&7\end{pmatrix},
\]
中心为$6$。从左上角开始顺时针读取邻点为$6,5,2,1,7,8,9,7$，比较结果为$1,0,0,0,1,1,1,1$，编码为$1+16+32+64+128=241$。把整幅邻域的灰度都乘以$2$再加$10$，编码不变。

图像旋转时，邻点在圆上循环移位，编码随之改变。Ojala等人对每个编码取所有循环移位中的最小值，得到旋转不变的编码；上例的二进制串$10001111$循环移位后的最小值是$00011111$，即$31$。他们进一步观察到，自然纹理中的大多数编码属于\term{均匀模式}（uniform pattern）：沿圆周$0$与$1$之间的跳变次数$U\leq2$，对应亮斑、暗斑、边缘和角点等基本结构。上例的跳变发生在$1\to0$与$0\to1$两处，$U=2$，是均匀模式。旋转不变的均匀模式编码$\mathrm{LBP}^{\mathrm{riu2}}_{P,R}$定义为：若$U\leq2$，取邻点中$1$的个数；否则统一取$P+1$。于是只剩$P+2$个取值，上例的编码为$5$。区域的纹理用这些编码的直方图描述，两区域之间用直方图距离比较。

\term{灰度共生矩阵}（gray-level co-occurrence matrix, GLCM）由Haralick等人提出\scite{vision-haralick1973}。先把灰度量化为$N_g$级，再固定一个位移$(\Delta x,\Delta y)$，统计区域内所有满足位移关系的像素对$(I(x,y),I(x+\Delta x,y+\Delta y))=(i,j)$出现的次数$C_{i,j}$，归一化为联合分布$p_{i,j}=C_{i,j}/\sum C$。常用位移是距离为1、方向为$0^\circ,45^\circ,90^\circ,135^\circ$的四种。由$p$计算的常用特征有
\begin{gather*}
  \text{对比度}=\sum_{i,j}(i-j)^2p_{i,j},\qquad
  \text{能量}=\sum_{i,j}p_{i,j}^2,\\
  \text{同质性}=\sum_{i,j}\frac{p_{i,j}}{1+\lvert i-j\rvert},\qquad
  \text{熵}=-\sum_{i,j}p_{i,j}\log_2p_{i,j}.
\end{gather*}
对比度大说明相邻像素灰度差异大，纹理粗糙；能量大说明少数灰度组合占主导，纹理规则。算例：$4\times4$图像、$N_g=3$，
\[
  \begin{pmatrix}0&0&1&1\\0&0&1&1\\0&2&2&2\\2&2&2&2\end{pmatrix},
\]
取水平位移$(1,0)$，每行有3对，共12对。逐行计数得$C_{0,0}=2$、$C_{0,1}=2$、$C_{0,2}=1$、$C_{1,1}=2$、$C_{2,2}=5$，其余为零。于是对比度$=(1\times2+4\times1)/12=0.5$，能量$=(4+4+1+4+25)/144\approx0.264$，同质性$=(2+2/2+1/3+2+5)/12\approx0.861$。若需要对方向不敏感，可以把$\bm C$与其转置相加，或对四个方向的特征取平均。

两种方法的差别在于统计的对象。LBP在每个像素处先做一次非线性的局部编码，再统计编码的直方图，对单调光照变化不变，计算快，适合逐像素给出纹理标签；GLCM直接统计灰度对的联合分布，特征的物理含义清楚，但依赖灰度量化级数与位移的选择，并且每种位移都要单独计算一个$N_g\times N_g$矩阵。二者都只描述二阶或局部的统计，不描述大尺度的纹理结构（如砖墙的排列周期）。卷积网络中，较浅层的特征图在空间上的统计也能刻画纹理，风格迁移就用卷积特征之间的格拉姆矩阵描述图像的纹理风格\scite{vision-gatys2016}，这相当于把“统计哪些量”也交给了学习。

表\ref{tab:vision-features-local-compare}比较本节的四种方法。

\begin{table*}[tbp]
  \centering
  \small
  \begin{tabularx}{\linewidth}{@{}lXXXX@{}}
    \toprule
    比较项 & Harris角点 & SIFT & HOG & LBP与GLCM \\
    \midrule
    取样方式 & 稀疏关键点 & 稀疏关键点（含尺度、方向） & 窗口内稠密网格 & 区域内稠密统计 \\
    输出 & 角点位置 & 位置、尺度、方向与128维描述子 & 3780维窗口描述子（$64\times128$窗口） & 编码直方图或标量纹理特征 \\
    统计的量 & 梯度外积的加权和 & 归一化邻域内的梯度方向直方图 & 单元内梯度方向直方图，块内归一化 & 邻点大小关系；灰度对共现频率 \\
    不变性 & 旋转、灰度偏移 & 尺度、旋转、仿射光照，容忍小形变 & 光照，容忍小形变；不处理旋转和尺度 & LBP：单调灰度变换；GLCM：依量化而定 \\
    典型用途 & 跟踪、几何估计 & 匹配、检索、拼接 & 滑动窗口检测 & 纹理分类与分割 \\
    主要失效 & 尺度变化 & 大视角、弱纹理、重复结构 & 姿态变化与遮挡 & 大尺度结构、量化选择 \\
    \bottomrule
  \end{tabularx}
  \caption{局部特征方法的比较。前两种回答“在哪里取特征、怎样跨图对应”，后两种回答“一个区域的整体形状或纹理是什么样”。}
  \label{tab:vision-features-local-compare}
\end{table*}

\section{区域与形状的描述}
\label{sec:vision-features-shape}

分割（下一节）或阈值化之后，我们得到的是一个个区域，例如路边交通标志的二值掩码。要判断它是圆形的禁令标志、三角形的警告标志还是八边形的停车标志，需要把区域的形状概括为少量数字，并且这些数字不应随区域在图中的位置、朝向和大小而变化，否则同一种标志在近处和远处会得到不同的描述。输入是一个区域（像素集合）或它的边界（有序像素序列），输出是一组形状描述量，要求的不变性一般包括平移、旋转和缩放，有时还包括描述的起点。

描述的对象有两种选择。边界描述只看轮廓线，把形状写成一条闭合曲线的编码或频谱，便于比较轮廓的细节，但对边界上的噪声和断裂敏感。区域描述使用区域内的全部像素，计算面积、矩等积分量，对边界噪声较稳健，却会把轮廓的细小凹凸平均掉。两种描述可以并用。

\subsection{边界的描述}
\label{sec:vision-features-boundary}

最直接的边界表示是沿轮廓逐步行走时记录方向。\term{链码}（chain code）是最常用的一种记录方式。Freeman链码规定8个方向，从$0$（向右）开始逆时针每$45^\circ$加一，于是$2$向上、$4$向左、$6$向下\scite{vision-freeman1961}。从轮廓上任一点出发，按$8$邻接（或$4$邻接，此时只用$0,2,4,6$）沿边界走一周，记下每一步的方向码，就得到这个闭合轮廓的链码。链码保存了轮廓的全部信息，并对平移不变，因为它只记录相邻点之间的相对位移。

链码依赖起点和朝向。例如沿一个宽2格、高1格的矩形边界逆时针行走，从左下角出发得到$0,0,2,4,4,6$。把图形旋转$90^\circ$，每一步方向都加$2$，链码变为$2,2,4,6,6,0$。为了消除旋转的影响，改用\term{一阶差分}（first difference）：记录相邻两步之间方向码的变化量$d_k=(c_k-c_{k-1})\bmod8$，首项与末项循环相减。两个链码的差分都是$2,0,2,2,0,2$，相同。为了消除起点的影响，把差分序列看成环，取所有循环移位中按数值最小的那一个，这里是$0,2,2,0,2,2$。这样得到的序列称为形状数，对平移、$90^\circ$整数倍的旋转和起点都不变。缩放不变需要先把边界重采样到固定步数，而离散网格上非$45^\circ$倍数的旋转会改变链码本身，所以链码适合规则网格上的形状比较，不宜直接用于任意角度旋转的比较。

傅里叶描述子从频率角度描述轮廓，可以较好地处理任意角度的旋转与缩放。把边界上等间隔采样的$K$个点写成复数$z_k=x_k+\mathrm iy_k$，$k=0,\ldots,K-1$，对这个序列做一维离散傅里叶变换
\begin{equation}
  a_u=\sum_{k=0}^{K-1}z_k\,\mathrm e^{-\mathrm i2\pi uk/K},\qquad u=0,\ldots,K-1.
  \label{eq:vision-features-fourier-descriptor}
\end{equation}
系数$a_u$称为\term{傅里叶描述子}（Fourier descriptor）。离散傅里叶变换是线性、可逆的，它的平移和调制性质在卷一已有讨论；把它们用到$z_k$上，就得到描述子在各种几何变换下的变化规律。平移$z_k\gets z_k+z_0$只改变$a_0$（增加$Kz_0$），其余系数不变；绕原点旋转$\phi$相当于$z_k\gets\mathrm e^{\mathrm i\phi}z_k$，所有系数乘以$\mathrm e^{\mathrm i\phi}$，模长不变；缩放$s$倍使所有系数乘以$s$；起点后移$k_0$个点使$a_u$乘以$\mathrm e^{\mathrm i2\pi uk_0/K}$，模长也不变。因此，舍弃$a_0$，取模长$\lvert a_u\rvert$，再除以$\lvert a_1\rvert$，得到的量对平移、旋转、缩放和起点都不变。

一个算例：边界只取四个点$(1,0),(0,1),(-1,0),(0,-1)$，即$z_k=\mathrm i^k=\mathrm e^{\mathrm i\pi k/2}$。代入式\eqref{eq:vision-features-fourier-descriptor}，$a_u=\sum_k\mathrm e^{\mathrm i2\pi k(1-u)/4}$，当$u=1$时每项都是$1$，$a_1=4$；其余$u$对应单位根的完整求和，结果为零。整个形状只由一个系数描述，这个系数对应以恒定角速度绕原点旋转一周的轨迹，也就是圆。一般地，低频系数描述轮廓的整体形状，高频系数描述细小的凹凸和锯齿；只保留前若干个系数再做逆变换，就得到轮廓的平滑近似。这一性质使傅里叶描述子能用很少的数字区分形状，也说明它对边界噪声的敏感度可以通过截断高频来控制。

两种边界描述都要求事先得到一条完整、闭合、有序的轮廓。区域若被遮挡或断成几段，链码与傅里叶描述子都会整体改变，因为它们是全局的：边界上一处的改变通过离散傅里叶变换影响所有系数。描述子之间的比较一般用欧氏距离，这隐含了两条轮廓的采样点一一对应的假设，对点数不同或部分缺失的轮廓并不合适。

\subsection{区域的描述}
\label{sec:vision-features-region}

区域描述从最简单的量开始。区域$R$的面积$A$是像素数，周长$P$是边界长度（链码中偶数码步长为$1$、奇数码步长为$\sqrt2$）。只靠面积或周长无法区分形状，因为它们都随大小变化；二者的组合
\begin{equation}
  \text{紧致度}=\frac{P^2}{4\pi A}
  \label{eq:vision-features-compactness}
\end{equation}
称为\term{紧致度}（compactness），它与大小无关：面积随尺度平方变化，周长的平方也随尺度平方变化，二者之比不变。由等周不等式，相同周长的平面图形中圆的面积最大，所以圆的紧致度最小，等于$1$；边长为$a$的正方形为$16a^2/(4\pi a^2)=4/\pi\approx1.27$；越细长或越曲折的形状，紧致度越大。离散网格上的周长估计有系统误差，同一个圆在不同分辨率下测出的紧致度会略有不同，比较时应使用相同的周长估计方式。

另一个描述拓扑的量是\term{欧拉数}（Euler number），定义为区域的连通分量数$C$减去孔洞数$H$：$E=C-H$。字母“B”是一个连通分量、两个孔，$E=-1$；“A”为$1-1=0$；“i”有两个分量、没有孔，$E=2$。欧拉数在拉伸、弯曲等不撕裂不粘连的形变下不变，适合区分拓扑结构不同的字符。它依赖于邻接的约定：前景用$8$邻接时，背景必须用$4$邻接（反之亦然），否则一条对角相接的像素链可能同时被看成“连通”和“围出一个孔”，这一对偶约定已在\ref{sec:vision-formation-connectivity}节说明。欧拉数在卷一中以单纯复形上的欧拉示性数出现，这里的定义是它在二值图像上的特例。

更一般的区域描述量是矩。把二值区域看成平面上的密度（区域内为$1$，区域外为$0$），它的$(p,q)$阶矩为
\[
  m_{p,q}=\sum_{(x,y)\in R}x^py^q.
\]
零阶矩$m_{0,0}$就是面积，一阶矩给出质心$\bar x=m_{1,0}/m_{0,0}$、$\bar y=m_{0,1}/m_{0,0}$。以质心为原点的\term{中心矩}（central moment）
\[
  \mu_{p,q}=\sum_{(x,y)\in R}(x-\bar x)^p(y-\bar y)^q
\]
对平移不变。二阶中心矩组成的矩阵$\begin{pmatrix}\mu_{2,0}&\mu_{1,1}\\\mu_{1,1}&\mu_{0,2}\end{pmatrix}$除以面积，就是区域像素坐标的协方差矩阵。它的特征向量给出区域的主轴方向，特征值之比给出伸长程度；主轴与$x$轴的夹角为$\theta=\tfrac12\operatorname{atan2}(2\mu_{1,1},\mu_{2,0}-\mu_{0,2})$。这与主成分分析在二维点集上的计算完全相同。

为了对缩放不变，Hu定义归一化中心矩
\begin{equation}
  \eta_{p,q}=\frac{\mu_{p,q}}{\mu_{0,0}^{1+(p+q)/2}}.
  \label{eq:vision-features-normalized-moment}
\end{equation}
指数的来历是：把区域放大$s$倍，求和中每个像素坐标乘以$s$，$(x-\bar x)^p(y-\bar y)^q$乘以$s^{p+q}$；连续情形下面积元也乘以$s^2$，于是$\mu_{p,q}$乘以$s^{p+q+2}$，而$\mu_{0,0}$乘以$s^2$。分母的指数$1+(p+q)/2$恰好使两者抵消。再对旋转不变，Hu由二阶、三阶归一化矩构造了七个不变量，称为\term{Hu矩}（Hu moments）\scite{vision-hu1962}，前两个是
\begin{equation}
  \phi_1=\eta_{2,0}+\eta_{0,2},\qquad
  \phi_2=(\eta_{2,0}-\eta_{0,2})^2+4\eta_{1,1}^2.
  \label{eq:vision-features-hu}
\end{equation}
$\phi_1$是归一化协方差矩阵的迹，$\phi_2$是两个特征值之差的平方；迹与特征值都不随坐标旋转改变，所以二者对旋转不变。其余五个由三阶矩组合而成，推导依据相同但代数较繁，这里不展开。

算例：一个宽4像素、高2像素的矩形区域，像素坐标$x\in\{0,1,2,3\}$、$y\in\{0,1\}$。面积$m_{0,0}=8$，质心$(1.5,0.5)$。$\mu_{2,0}=2\times(2.25+0.25+0.25+2.25)=10$，$\mu_{0,2}=8\times0.25=2$，$\mu_{1,1}=0$。于是$\eta_{2,0}=10/64\approx0.156$，$\eta_{0,2}=2/64\approx0.031$，$\phi_1\approx0.188$，$\phi_2=(8/64)^2\approx0.016$。把矩形旋转$90^\circ$，$\eta_{2,0}$与$\eta_{0,2}$互换，$\phi_1,\phi_2$不变。把矩形放大为$8\times4$像素，同样计算得$\mu_{2,0}=168$、$\mu_{0,2}=40$、$m_{0,0}=32$，$\eta_{2,0}\approx0.164$、$\eta_{0,2}\approx0.039$，$\phi_1\approx0.203$，与$0.188$并不相等。差别来自离散化：连续矩形的$\eta_{2,0}$等于宽高比的$1/12$倍，$4\times2$与$8\times4$都应为$1/6\approx0.167$；离散求和把每个像素当成一个点，漏掉了像素自身面积的二阶矩（每个单位像素为$1/12$）。在$4\times2$的情形补上这一项，$\mu_{2,0}=10+8/12$，$\eta_{2,0}=0.167$，与连续值一致。这说明Hu矩的不变性是连续意义上的，在像素很少的小区域上会有明显的离散误差。

区域描述量计算简单、意义直观，适合作为分类器的输入特征或分割结果的筛选条件，例如按面积去掉过小的连通分量、按紧致度保留近似圆形的标志候选。它们的局限在于判别力有限：七个Hu矩远不足以区分形状相近的类别，高阶矩对噪声又很敏感，因为$x^py^q$在远离质心处被放大。这些描述量都假设区域已经被正确分割出来；若分割把标志与旁边的杆子连成一块，任何形状描述都无济于事。形状描述因此依赖于下一节的分割质量。学习式方法通常不再显式计算这些量，而是由网络直接从像素预测类别和掩码，但在实例分割之后，仍然常用面积、外接框和连通分量对结果做后处理。

\section{经典图像分割}
\label{sec:vision-features-segmentation}

上一节假定区域已经存在。现在回到这个前提：怎样把街景中的像素分成路面、天空、车辆、招牌等区域？\term{图像分割}（image segmentation）把图像的像素集合$\Omega$划分为互不相交的区域$R_1,\ldots,R_n$，满足$\bigcup_iR_i=\Omega$；每个区域内部是连通的，并且满足某个一致性判据$Q(R_i)$，例如灰度方差小于阈值；相邻区域的并集不满足这一判据。这个定义说明分割同时受两种约束：区域内部要相似，区域在空间上要成片。经典分割不知道“路面”“车辆”是什么，它只根据灰度、颜色、纹理和位置把像素组织起来；给区域赋予类别是\ref{chap:vision-image-recognition}章语义分割的任务。

难点在于“相似”与“成片”常常冲突。同一辆车的车身有高光和阴影，灰度并不一致；车顶与路面的颜色可能几乎相同，边界却是真实存在的；光照从画面左侧到右侧逐渐变暗，使同一路面的灰度在空间上缓慢变化。不同分割方法的区别，在于使用哪些信息、在多大范围内做决策。

四条路线依次扩大决策所依据的信息。基于阈值的分割只看每个像素的灰度，用直方图选一个或多个阈值，速度最快，但完全不用空间关系。基于区域的分割从种子或整幅图出发，在空间邻接的约束下合并相似像素或拆分不一致的块；分水岭把梯度幅值看成地形，用浸没过程找出被梯度“山脊”隔开的盆地。基于聚类的分割把每个像素表示为颜色与位置组成的特征向量，在特征空间中分组，SLIC超像素用局部K均值得到大小均匀、贴合边界的小区域。基于图的分割把像素看成图的顶点、相似度看成边权，把分割写成图上的全局优化问题：归一化割求平衡的二划分，图割在用户给出的前景与背景约束下求能量的全局最小，GrabCut再交替估计颜色模型与分割。

\subsection{基于阈值的分割}
\label{sec:vision-features-otsu}
\label{sec:vision-features-threshold}

若目标与背景的灰度分布不重叠，例如白纸上的黑字、夜空下的路灯，只需一个阈值$T$：$I(x,y)>T$的像素归为一类，其余归为另一类。阈值法的核心问题是怎样从图像本身自动选出$T$。灰度直方图上若有两个峰，阈值应落在两峰之间的谷底；但噪声使谷底不明显，直方图又常有许多小的起伏，直接找谷底并不可靠。

Otsu把阈值选择写成一个判别准则\scite{vision-otsu1979}：好的阈值应当使两类内部的灰度尽量集中、两类之间的灰度尽量分开。设灰度级为$0,\ldots,L-1$，归一化直方图为$p_i$。阈值$k$把灰度级分为$C_0=\{0,\ldots,k\}$与$C_1=\{k+1,\ldots,L-1\}$。记
\[
  \omega(k)=\sum_{i=0}^{k}p_i,\qquad \mu(k)=\sum_{i=0}^{k}ip_i,\qquad \mu_T=\sum_{i=0}^{L-1}ip_i,
\]
则两类的概率为$\omega_0=\omega(k)$、$\omega_1=1-\omega(k)$，均值为$\mu_0=\mu(k)/\omega(k)$、$\mu_1=(\mu_T-\mu(k))/(1-\omega(k))$，并满足$\omega_0\mu_0+\omega_1\mu_1=\mu_T$。

第一步说明类内方差与类间方差之和是常数。总方差$\sigma_T^2=\sum_i(i-\mu_T)^2p_i$与$k$无关。把求和按两类拆开，在每类中把$i-\mu_T$写成$(i-\mu_c)+(\mu_c-\mu_T)$并展开平方。交叉项为$2(\mu_c-\mu_T)\sum_{i\in C_c}(i-\mu_c)p_i$，由均值的定义，这个和为零。于是
\[
  \sigma_T^2=\underbrace{\omega_0\sigma_0^2+\omega_1\sigma_1^2}_{\sigma_W^2}
  +\underbrace{\omega_0(\mu_0-\mu_T)^2+\omega_1(\mu_1-\mu_T)^2}_{\sigma_B^2},
\]
其中$\sigma_c^2$是第$c$类的类内方差。这就是全方差公式在两类情形下的形式。由于$\sigma_T^2$固定，最小化类内方差$\sigma_W^2$等价于最大化类间方差$\sigma_B^2$，而后者只需要一阶统计量。

第二步化简$\sigma_B^2$。由$\mu_T=\omega_0\mu_0+\omega_1\mu_1$得$\mu_0-\mu_T=\omega_1(\mu_0-\mu_1)$、$\mu_1-\mu_T=\omega_0(\mu_1-\mu_0)$，代入得$\sigma_B^2=\omega_0\omega_1^2(\mu_1-\mu_0)^2+\omega_1\omega_0^2(\mu_1-\mu_0)^2=\omega_0\omega_1(\mu_1-\mu_0)^2$。再把$\mu_0,\mu_1$的表达式代入，$\mu_1-\mu_0=(\omega\mu_T-\mu)/(\omega(1-\omega))$（这里$\omega,\mu$是$\omega(k),\mu(k)$的简写），于是
\begin{equation}
  \sigma_B^2(k)=\frac{\bigl(\mu_T\,\omega(k)-\mu(k)\bigr)^2}{\omega(k)\bigl(1-\omega(k)\bigr)},\qquad
  k^*=\argmax_{k:\,0<\omega(k)<1}\sigma_B^2(k).
  \label{eq:vision-features-otsu}
\end{equation}
$\omega(k)$与$\mu(k)$都是累加量，随$k$递推更新，所以一次扫描直方图就能算出所有候选阈值的$\sigma_B^2$。比值$\sigma_B^2(k^*)/\sigma_T^2\in[0,1]$衡量两类分开的程度，可以用来判断这幅图是否适合用单一阈值分割。

\begin{algorithm}[Otsu阈值选择]
  \label{alg:vision-features-otsu}
  \begin{algorithmic}[1]
    \Require 灰度直方图计数$h_0,\ldots,h_{L-1}$
    \Ensure 阈值$k^*$与可分性$\sigma_B^2(k^*)/\sigma_T^2$
    \State $N\gets\sum_ih_i$，$p_i\gets h_i/N$，$\mu_T\gets\sum_iip_i$，$\sigma_T^2\gets\sum_i(i-\mu_T)^2p_i$
    \State $\omega\gets0$，$\mu\gets0$
    \For{$k=0,\ldots,L-2$}
      \State $\omega\gets\omega+p_k$，$\mu\gets\mu+kp_k$
      \If{$0<\omega<1$}
        \State 按式\eqref{eq:vision-features-otsu}计算$\sigma_B^2(k)$，记录最大者
      \EndIf
    \EndFor
    \State 返回取得最大值的$k^*$及$\sigma_B^2(k^*)/\sigma_T^2$
  \end{algorithmic}
\end{algorithm}

算例：一幅40个像素、8级灰度的教学图像，直方图计数为$(2,\allowbreak6,\allowbreak9,\allowbreak5,\allowbreak2,\allowbreak4,\allowbreak8,\allowbreak4)$。总均值$\mu_T=143/40=3.575$，总方差$\sigma_T^2\approx4.79$。表\ref{tab:vision-features-otsu-example}列出各候选阈值的累加量与类间方差。

\begin{table}[htbp]
  \centering
  \small
  \begin{tabular}{@{}cccccc@{}}
    \toprule
    $k$ & $\omega(k)$ & $\mu(k)$ & $\mu_0$ & $\mu_1$ & $\sigma_B^2(k)$ \\
    \midrule
    0 & 0.050 & 0.000 & 0.000 & 3.763 & 0.67 \\
    1 & 0.200 & 0.150 & 0.750 & 4.281 & 2.00 \\
    2 & 0.425 & 0.600 & 1.412 & 5.174 & 3.46 \\
    3 & 0.550 & 0.975 & 1.773 & 5.778 & 3.97 \\
    4 & 0.600 & 1.175 & 1.958 & 6.000 & 3.92 \\
    5 & 0.700 & 1.675 & 2.393 & 6.333 & 3.26 \\
    6 & 0.900 & 2.875 & 3.194 & 7.000 & 1.30 \\
    \bottomrule
  \end{tabular}
  \caption{Otsu算例中各候选阈值的统计量。$\mu_0,\mu_1$由$\omega(k),\mu(k)$与$\mu_T=3.575$算出，$\sigma_B^2$按式\eqref{eq:vision-features-otsu}计算后保留两位小数。数据为教学构造。}
  \label{tab:vision-features-otsu-example}
\end{table}

以$k=3$为例复核：$\mu(3)=(0\times2+1\times6+2\times9+3\times5)/40=0.975$，$\sigma_B^2=(3.575\times0.55-0.975)^2/(0.55\times0.45)=0.98125^2/0.2475\approx3.970$；也可以用$\omega_0\omega_1(\mu_1-\mu_0)^2=0.2475\times4.005^2\approx3.970$核对。最大值在$k^*=3$处取得，灰度$0$至$3$为一类、$4$至$7$为另一类，可分性约为$3.97/4.79\approx0.83$。图\ref{fig:vision-features-otsu}画出了直方图与$\sigma_B^2(k)$曲线。注意$k=3$与$k=4$的类间方差只差$0.05$：灰度$4$只有2个像素，放在哪一类对结果影响很小，这正对应直方图在两峰之间的低谷。

\begin{figure}[htbp]
  \centering
  % Otsu 教学例：40 个像素、8 级灰度，计数 2,6,9,5,2,4,8,4。上：直方图与选出的阈值
% k*=3；下：各候选阈值的类间方差 sigma_B^2(k)（由正文公式精确计算后保留两位）。宽约 84mm。
\begin{tikzpicture}[x=9mm,y=2.4mm,font=\figurefont\small]
  \foreach \y in {3,6,9}{\draw[FigureGrid,line width=.4pt] (-.6,\y)--(7.6,\y);
    \node[anchor=east,font=\figurefont\footnotesize] at (-.65,\y) {\y};}
  \foreach \v [count=\k from 0] in {2,6,9,5,2,4,8,4}{
    \ifnum\k<4
      \fill[FigureBlueFill,draw=FigureBlue,line width=.6pt] (\k-.35,0) rectangle (\k+.35,\v);
    \else
      \fill[FigureYellowFill,draw=FigureYellow,line width=.6pt] (\k-.35,0) rectangle (\k+.35,\v);
    \fi
    \node[below,font=\figurefont\footnotesize] at (\k,0) {\k};
  }
  \draw[FigureStroke,line width=.6pt] (-.6,0)--(7.6,0);
  \draw[FigureRed,dashed,line width=.8pt] (3.5,-.5)--(3.5,11);
  \node[anchor=south,font=\figurefont\footnotesize] at (3.5,11) {$k^*=3$};
  \node[font=\figurefont\footnotesize,text=FigureMuted] at (1.5,11.5) {$C_0$};
  \node[font=\figurefont\footnotesize,text=FigureMuted] at (5.5,11.5) {$C_1$};
  \node[anchor=east] at (-1.4,5) {像素数};
  \node at (3.5,-3.4) {灰度级 $i$};
  % sigma_B^2 curve
  \begin{scope}[yshift=-12mm,y=2.4mm]
    \draw[FigureGrid,line width=.4pt] (-.6,0)--(7.6,0);
    \foreach \y in {2,4}{\draw[FigureGrid,line width=.4pt] (-.6,{-6+\y})--(7.6,{-6+\y});
      \node[anchor=east,font=\figurefont\footnotesize] at (-.65,{-6+\y}) {\y};}
    \draw[FigureStroke,line width=.6pt] (-.6,-6)--(7.6,-6);
    \draw[FigureBlue,line width=1pt] (0,{-6+0.67})--(1,{-6+2.00})--(2,{-6+3.46})--(3,{-6+3.97})--(4,{-6+3.92})--(5,{-6+3.26})--(6,{-6+1.30});
    \foreach \k/\v in {0/0.67,1/2.00,2/3.46,3/3.97,4/3.92,5/3.26,6/1.30}{
      \fill[FigureBlue] (\k,{-6+\v}) circle (.7mm);}
    \node[anchor=west,font=\figurefont\footnotesize] at (3.1,{-6+4.5}) {$3.97$};
    \foreach \k in {0,...,6}{\node[below,font=\figurefont\footnotesize] at (\k,-6) {\k};}
    \node[anchor=east] at (-1.4,-3) {$\sigma_B^2(k)$};
    \node at (3.5,-9.4) {候选阈值 $k$};
  \end{scope}
\end{tikzpicture}

  \caption{Otsu阈值选择的教学例。上：40个像素的8级灰度直方图，蓝色与黄色柱分别属于阈值$k^*=3$分出的两类$C_0,C_1$。下：各候选阈值的类间方差$\sigma_B^2(k)$，最大值$3.97$对应选出的阈值。数据与表\ref{tab:vision-features-otsu-example}一致。}
  \label{fig:vision-features-otsu}
\end{figure}

Otsu方法在两类灰度分布近似单峰、方差相近、比例不过分悬殊时效果很好。它的失效情形可以直接从准则看出。准则中含有因子$\omega_0\omega_1$，它倾向于给出两类概率不太悬殊的划分；若前景只占极少像素（例如大片路面上的一小块标志），直方图上前景的峰很低，选出的阈值常常偏离两峰之间的谷底，落进背景分布的内部；若光照不均，同一类物体在图像不同位置的灰度不同，全局直方图上两类相互重叠，没有任何单一阈值能把它们分开。

后一种情形用\term{自适应阈值}（adaptive thresholding）处理：每个像素的阈值由它周围窗口内的统计量决定。最简单的形式是$T(x,y)=m(x,y)-C$，$m$为窗口均值、$C$为小的常数；招牌上的文字在阴影中也比周围的底色暗，相对局部均值仍能被分出。Sauvola与Pietikäinen为文档图像提出$T(x,y)=m(x,y)\bigl[1+\kappa\bigl(s(x,y)/R-1\bigr)\bigr]$，其中$s$为窗口标准差、$R$为标准差的动态范围、$\kappa$为正的小常数\scite{vision-sauvola2000}：在对比度高（$s$大）的文字区域，阈值接近局部均值；在对比度低的平坦背景，阈值降到均值以下，减少背景被误判为前景。窗口必须大于文字笔画等前景结构，否则窗口内全是前景，局部均值也就失去了参照作用。逐像素计算窗口均值和方差看起来代价很高，但用\ref{sec:vision-features-viola-jones}节的积分图，任意窗口的像素和与平方和都能用四次查表得到，每个像素的计算量与窗口大小无关。

\subsection{基于区域的分割}
\label{sec:vision-features-region-segmentation}
\label{sec:vision-features-watershed}

阈值法逐个像素独立决策，结果中同一类的像素可能散布在图像各处，也可能被噪声打成碎点。基于区域的方法把空间连通性直接写进算法：区域只通过邻接关系扩展或合并。这一类方法有三种代表。区域生长从种子出发向外扩张；分裂合并从整幅图出发向下分解再向上合并；分水岭不使用一致性判据，而是利用梯度幅值的地形，让区域在梯度的“山脊”处相遇。

\term{区域生长}（region growing）从一个或多个种子像素开始，反复检查当前区域边界上的邻接像素，若满足相似性判据就加入区域，直到没有像素可以加入。常用判据是$\lvert I(p)-\bar I_R\rvert\leq\tau$，$\bar I_R$为当前区域的均值。算例：$4\times4$图像
\[
  \begin{pmatrix}10&11&30&31\\12&10&29&33\\11&25&28&30\\40&42&41&39\end{pmatrix},
\]
取左上角为种子、$\tau=5$、4邻接。区域依次吸收$11$、$12$、$10$、$11$，均值变为$10.8$；右侧的$30,29$与下方的$25,40$都与均值相差超过$5$，生长停止，得到左上角的5个像素。若在右上角$31$处再放一个种子，可以得到灰度在$28$至$33$之间的6个像素；中间的$25$与两个区域的均值相差都超过$5$，不属于任何区域，下面一行则需要第三个种子。这个例子显示了区域生长的两个特点：结果取决于种子的位置与个数，以及判据的阈值；处在两区域之间的过渡像素可能无处可去。此外，以区域均值为参照时，加入像素的顺序会改变均值，从而影响后续判断，结果对遍历顺序也有一定依赖。

\term{分裂合并}（split and merge）不需要种子。从整幅图开始，若某个块不满足一致性判据$Q$（例如块内最大灰度与最小灰度之差超过$\tau$），就把它四等分，递归进行，得到一棵四叉树；然后检查相邻的叶块，若两块之并满足$Q$就合并。分裂保证每个块内部一致，合并去掉四叉树划分引入的人为边界。它的缺点是块的形状受四叉树约束，边界呈阶梯状，需要后续平滑。

\term{分水岭}（watershed）变换换了一种看待图像的方式。把梯度幅值图看成地形：均匀区域内部梯度小，是低洼的盆地；区域边界处梯度大，是山脊。设想从每个局部极小处向上注水，水位逐渐上升，每个盆地的水面各自扩大；当两个盆地的水即将相遇时，在相遇处筑一道坝。水位升到最高时，所有的坝就构成了区域边界，每个盆地就是一个区域。这样得到的边界一定位于梯度较大的位置，并且自然是闭合的。

Vincent与Soille给出了这一浸没过程的高效离散实现\scite{vision-vincent1991}：先把像素按灰度排序，再按灰度从低到高逐层处理，每一层用先进先出队列从已标记的盆地向外扩展，新出现的、不与任何已有盆地相连的像素成为新的极小，从而开始一个新盆地。下面的算法采用等价的优先队列形式，并直接从标记出发，这也是后面讲的标记控制分水岭的通常写法。

\begin{algorithm}[基于优先队列的分水岭]
  \label{alg:vision-features-watershed}
  \begin{algorithmic}[1]
    \Require 地形图$f$（通常为梯度幅值），标记图$\ell$（若干连通的已标记区域，其余为未标记）
    \Ensure 每个像素的区域标签，坝像素标记为边界
    \State 把所有与已标记区域相邻的未标记像素$p$以优先级$f(p)$放入队列
    \While{队列非空}
      \State 取出优先级最小的像素$p$
      \If{$p$的已标记邻点都属于同一个区域$r$}
        \State 令$\ell(p)\gets r$，把$p$的未入队未标记邻点$q$以优先级$f(q)$放入队列
      \Else
        \State 把$p$标记为坝 \Comment{两个区域在此相遇}
      \EndIf
    \EndWhile
  \end{algorithmic}
\end{algorithm}

以一维剖面为例（图\ref{fig:vision-features-watershed}）。设11个位置上的梯度幅值为
\[
  f=(3,1,2,5,4,0,2,6,3,1,3).
\]
局部极小在位置$1,5,9$，取它们作为标记。按算法从低到高扩展：位置$2$（值$2$）、$6$（值$2$）先被各自相邻的盆地吸收，随后是值为$3$的位置$0,8,10$，再是值$4$的位置$4$并入中间盆地。位置$3$（值$5$）的两侧分别属于左盆地和中盆地，成为坝；位置$7$（值$6$）同理。结果是三个区域$\{0,1,2\}$、$\{4,5,6\}$、$\{8,9,10\}$。

\begin{figure}[htbp]
  \centering
  % 分水岭的一维教学例：梯度幅值剖面 f=(3,1,2,5,4,0,2,6,3,1,3)，位置 0 至 10。
% 上：三个局部极小各自形成盆地，坝在位置 3 与 7；下：只在位置 1、9 设标记时，
% 中间盆地经较低的鞍点（位置 3，高 5）并入左侧区域，坝只剩位置 7。宽约 96mm。
\begin{tikzpicture}[x=8mm,y=2.1mm,font=\figurefont\small]
\begin{scope}[yshift=0mm]
\fill[FigureBlueFill,draw=FigureBlue,line width=.6pt] (-0.42,0) rectangle (0.42,3);
\node[font=\figurefont\footnotesize] at (0,3.9) {A};
\fill[FigureBlueFill,draw=FigureBlue,line width=.6pt] (0.58,0) rectangle (1.42,1);
\node[font=\figurefont\footnotesize] at (1,1.9) {A};
\fill[FigureBlueFill,draw=FigureBlue,line width=.6pt] (1.58,0) rectangle (2.42,2);
\node[font=\figurefont\footnotesize] at (2,2.9) {A};
\fill[FigureStroke] (2.58,0) rectangle (3.42,5);
\fill[FigureYellowFill,draw=FigureYellow,line width=.6pt] (3.58,0) rectangle (4.42,4);
\node[font=\figurefont\footnotesize] at (4,4.9) {B};
\fill[FigureYellowFill,draw=FigureYellow,line width=.6pt] (4.58,0) rectangle (5.42,0);
\node[font=\figurefont\footnotesize] at (5,0.9) {B};
\fill[FigureYellowFill,draw=FigureYellow,line width=.6pt] (5.58,0) rectangle (6.42,2);
\node[font=\figurefont\footnotesize] at (6,2.9) {B};
\fill[FigureStroke] (6.58,0) rectangle (7.42,6);
\fill[FigureRedFill,draw=FigureRed,line width=.6pt] (7.58,0) rectangle (8.42,3);
\node[font=\figurefont\footnotesize] at (8,3.9) {C};
\fill[FigureRedFill,draw=FigureRed,line width=.6pt] (8.58,0) rectangle (9.42,1);
\node[font=\figurefont\footnotesize] at (9,1.9) {C};
\fill[FigureRedFill,draw=FigureRed,line width=.6pt] (9.58,0) rectangle (10.42,3);
\node[font=\figurefont\footnotesize] at (10,3.9) {C};
\draw[FigureStroke,line width=.7pt] (1,-1.2)--(1,-0.2);
\fill[FigureStroke] (1,-1.6) circle (0.7mm);
\draw[FigureStroke,line width=.7pt] (5,-1.2)--(5,-0.2);
\fill[FigureStroke] (5,-1.6) circle (0.7mm);
\draw[FigureStroke,line width=.7pt] (9,-1.2)--(9,-0.2);
\fill[FigureStroke] (9,-1.6) circle (0.7mm);
\draw[FigureStroke,line width=.6pt] (-.6,0)--(10.6,0);
\node[anchor=west] at (-.6,8) {(a) 无标记的浸没：三个盆地};
\end{scope}
\begin{scope}[yshift=-27.3mm]
\fill[FigureBlueFill,draw=FigureBlue,line width=.6pt] (-0.42,0) rectangle (0.42,3);
\node[font=\figurefont\footnotesize] at (0,3.9) {A};
\fill[FigureBlueFill,draw=FigureBlue,line width=.6pt] (0.58,0) rectangle (1.42,1);
\node[font=\figurefont\footnotesize] at (1,1.9) {A};
\fill[FigureBlueFill,draw=FigureBlue,line width=.6pt] (1.58,0) rectangle (2.42,2);
\node[font=\figurefont\footnotesize] at (2,2.9) {A};
\fill[FigureBlueFill,draw=FigureBlue,line width=.6pt] (2.58,0) rectangle (3.42,5);
\node[font=\figurefont\footnotesize] at (3,5.9) {A};
\fill[FigureBlueFill,draw=FigureBlue,line width=.6pt] (3.58,0) rectangle (4.42,4);
\node[font=\figurefont\footnotesize] at (4,4.9) {A};
\fill[FigureBlueFill,draw=FigureBlue,line width=.6pt] (4.58,0) rectangle (5.42,0);
\node[font=\figurefont\footnotesize] at (5,0.9) {A};
\fill[FigureBlueFill,draw=FigureBlue,line width=.6pt] (5.58,0) rectangle (6.42,2);
\node[font=\figurefont\footnotesize] at (6,2.9) {A};
\fill[FigureStroke] (6.58,0) rectangle (7.42,6);
\fill[FigureRedFill,draw=FigureRed,line width=.6pt] (7.58,0) rectangle (8.42,3);
\node[font=\figurefont\footnotesize] at (8,3.9) {C};
\fill[FigureRedFill,draw=FigureRed,line width=.6pt] (8.58,0) rectangle (9.42,1);
\node[font=\figurefont\footnotesize] at (9,1.9) {C};
\fill[FigureRedFill,draw=FigureRed,line width=.6pt] (9.58,0) rectangle (10.42,3);
\node[font=\figurefont\footnotesize] at (10,3.9) {C};
\draw[FigureStroke,line width=.7pt] (1,-1.2)--(1,-0.2);
\fill[FigureStroke] (1,-1.6) circle (0.7mm);
\draw[FigureStroke,line width=.7pt] (9,-1.2)--(9,-0.2);
\fill[FigureStroke] (9,-1.6) circle (0.7mm);
\draw[FigureStroke,line width=.6pt] (-.6,0)--(10.6,0);
\node[anchor=west] at (-.6,8) {(b) 以位置 1、9 为标记：两个区域};
\end{scope}
\node[font=\figurefont\footnotesize] at (0,-16.2) {0};
\node[font=\figurefont\footnotesize] at (1,-16.2) {1};
\node[font=\figurefont\footnotesize] at (2,-16.2) {2};
\node[font=\figurefont\footnotesize] at (3,-16.2) {3};
\node[font=\figurefont\footnotesize] at (4,-16.2) {4};
\node[font=\figurefont\footnotesize] at (5,-16.2) {5};
\node[font=\figurefont\footnotesize] at (6,-16.2) {6};
\node[font=\figurefont\footnotesize] at (7,-16.2) {7};
\node[font=\figurefont\footnotesize] at (8,-16.2) {8};
\node[font=\figurefont\footnotesize] at (9,-16.2) {9};
\node[font=\figurefont\footnotesize] at (10,-16.2) {10};
\end{tikzpicture}

  \caption{分水岭的一维教学例。柱高为梯度幅值，字母为所属区域，深色柱为坝，柱下圆点为极小或标记。(a) 以全部三个局部极小为种子，得到三个盆地，坝在位置$3$与$7$。(b) 只在位置$1$与$9$放标记，位置$5$的极小不再单独成区：水从左侧经高度为$5$的位置$3$先到达它，所以它并入左侧区域，坝只剩高度为$6$的位置$7$。数值为教学构造。}
  \label{fig:vision-features-watershed}
\end{figure}

直接对梯度图做分水岭，几乎总会把图像切得过碎：噪声和纹理在梯度图中产生大量浅的局部极小，每个都成为一个小区域，一辆车可能被切成上百块。这种现象称为\term{过分割}（over-segmentation）。解决办法是\term{标记控制分水岭}（marker-controlled watershed）：不让每个局部极小都成为盆地，而是事先指定每个目标与背景各一个标记，只从这些标记开始注水。图\ref{fig:vision-features-watershed}(b)演示了这一点：只保留位置$1$与$9$两个标记时，位置$5$的浅盆地没有自己的标记，水从左侧经位置$3$（高度$5$）先于从右侧经位置$7$（高度$6$）到达它，于是它并入左侧区域。标记可以由用户点选，也可以自动产生，例如取深度超过$h$的极小（用\ref{sec:vision-transforms-morphology-algorithms}节的形态学重建去掉浅于$h$的极小），或取距离变换的极大值作为互相粘连的圆形物体的中心。标记决定了区域的个数和大致位置，分水岭只负责把边界放到梯度最大处。

三种区域方法都只利用局部的一致性或梯度，不知道区域应该有多大、应该是什么形状。区域生长与分裂合并的结果依赖阈值，分水岭的结果依赖标记；当目标内部的变化与目标边界的对比度相当时，三者都会出错。

\subsection{基于聚类的分割}
\label{sec:vision-features-clustering-segmentation}
\label{sec:vision-features-superpixel}

把每个像素$p$表示为特征向量$\bm f_p$，例如Lab颜色$(L,a,b)$，或者再拼接上乘以权重的位置$(\lambda x,\lambda y)$，分割就变成了对这些向量的聚类。Lab空间的欧氏距离比RGB空间更接近感知上的颜色差，这一点见\ref{sec:vision-formation-color-spaces}节。聚类方法本身在卷四已经讲过，这里说明它们用于图像时的特殊之处。

K均值把$N$个向量分成$K$组，交替执行两步：把每个向量分给最近的中心；把每个中心更新为所分向量的均值。每一步都不增加组内平方距离之和，因而算法在有限步内收敛到局部最优，但结果依赖初始中心，可参考\BookChapterRef{chap:clustering}{卷四“聚类”章}。只用颜色时，K均值得到的“区域”在空间上并不连通，例如红色车身和远处的红色招牌会被分到同一组；加入位置后，区域在空间上变得紧凑，但位置权重$\lambda$需要人工选择。$K$也要预先给定，而一幅街景中应有多少个区域并不知道。

均值漂移（mean shift）不需要预先给定组数。它把特征向量看成某个概率密度的样本，从每个样本出发，沿核密度估计的梯度方向反复移动到局部窗口内样本的加权均值，直到收敛到密度的一个众数；收敛到同一众数的像素归为一组。均值漂移的收敛性与带宽选择也在卷四同一章讨论过。Comaniciu与Meer把它用于图像时，在位置与颜色组成的联合空间中使用乘积核，分别用空间带宽$h_s$和颜色带宽$h_r$控制两种距离的尺度\scite{clust-comaniciu2002}：只有在空间上相距约$h_s$以内、颜色上相差约$h_r$以内的像素才会相互吸引。它能保持边界、不需要指定组数，代价是计算量大，并且$h_s,h_r$同样需要按图像调整。

很多情况下我们并不要求一次就分出完整的物体，只需要把图像压缩为几百个大小相近、内部均匀、边界贴合物体轮廓的小区域，作为后续处理的基本单位。这样的区域称为\term{超像素}（superpixel）。用超像素代替像素，后续的图优化或分类的变量数可以减少两到三个数量级，而物体边界基本得以保留。

SLIC（simple linear iterative clustering）在Lab颜色与位置组成的五维空间中做K均值，但只在局部搜索\scite{vision-achanta2012}。设图像有$N$个像素，希望得到$K$个超像素，则网格间距$S=\sqrt{N/K}$。距离度量为
\begin{equation}
  D=\sqrt{d_c^2+\Bigl(\frac{d_s}{S}\Bigr)^2m^2},
  \label{eq:vision-features-slic}
\end{equation}
其中$d_c$为Lab颜色的欧氏距离，$d_s$为像素坐标的欧氏距离。颜色差与空间距离的量纲不同，$d_s$除以$S$把空间距离归一化到“几个网格间距”，再用紧致度参数$m$调节两者的相对权重。$m$大时空间距离占主导，超像素形状规则；$m$小时颜色占主导，超像素更贴合边界但形状不规则。原文指出在CIELAB空间中$m$可取$[1,40]$，默认值为$10$。

\begin{algorithm}[SLIC超像素]
  \label{alg:vision-features-slic}
  \begin{algorithmic}[1]
    \Require Lab图像，超像素个数$K$，紧致度$m$，残差阈值
    \Ensure 每个像素的超像素标签
    \State 以间距$S=\sqrt{N/K}$在规则网格上放置中心$\bm C_k=(l_k,a_k,b_k,x_k,y_k)$
    \State 把每个中心移到其$3\times3$邻域中梯度最小的位置 \Comment{避免落在边缘或噪声点上}
    \Repeat
      \State 对所有像素置距离$d(i)\gets\infty$
      \For{每个中心$\bm C_k$}
        \For{$\bm C_k$周围$2S\times2S$区域内的每个像素$i$}
          \State 按式\eqref{eq:vision-features-slic}计算$D$；若$D<d(i)$，则$d(i)\gets D$、$\ell(i)\gets k$
        \EndFor
      \EndFor
      \State 把每个中心更新为所属像素的均值，计算中心移动的残差$E$
    \Until{$E$小于阈值}
    \State 把不与所属中心连通的孤立像素块并入相邻的超像素
  \end{algorithmic}
\end{algorithm}

与标准K均值相比，SLIC的关键改动是把每个中心的搜索范围限制在$2S\times2S$的窗口内。超像素的期望大小是$S\times S$，所以这个窗口足以覆盖它可能的范围；每个像素只与附近少数中心比较，每次迭代的计算量与像素数$N$成正比，而与$K$无关。

算例：取$S=10$，某像素的Lab值为$(50,10,10)$、坐标为$(12,8)$。中心A为$(52,12,9)$、位于$(5,5)$；中心B为$(60,10,10)$、位于$(14,9)$。到A的颜色距离$d_c=\sqrt{4+4+1}=3$，空间距离$d_s=\sqrt{49+9}\approx7.62$；到B的颜色距离$d_c=10$，空间距离$d_s=\sqrt{4+1}\approx2.24$。取$m=10$，$D_A=\sqrt{9+0.581\times100}\approx8.19$，$D_B=\sqrt{100+0.05\times100}\approx10.25$，像素归A，颜色相近起了决定作用；取$m=40$，$D_A=\sqrt{9+0.581\times1600}\approx30.6$，$D_B=\sqrt{100+0.05\times1600}\approx13.4$，像素归B，空间相近起了决定作用。

聚类方法在特征空间中决策，只在特征里包含位置时才间接使用空间关系，并且距离是逐对计算的，区域之间的边界由“离哪个中心近”决定，而不是由边界本身的强弱决定。超像素把这一问题留给下一步：它只负责过分割，合并成物体还需要另外的判据。

\subsection{基于图的分割}
\label{sec:vision-features-graph-segmentation}

图的方法把像素之间的关系显式写出来。把每个像素（或超像素）看成顶点，相邻或相近的顶点之间连一条边，边权$w_{i,j}\geq0$表示两点属于同一区域的倾向，例如
\[
  w_{i,j}=\exp\Bigl(-\frac{\lVert\bm f_i-\bm f_j\rVert_2^2}{\sigma_f^2}\Bigr)\exp\Bigl(-\frac{\lVert\bm x_i-\bm x_j\rVert_2^2}{\sigma_x^2}\Bigr),\quad \lVert\bm x_i-\bm x_j\rVert_2<r,
\]
其中$\bm f$为颜色或亮度，$\bm x$为位置，距离超过$r$的点对不连边。分割就是把顶点分组，使组内的边权大、组间被切断的边权小。两种代表方法对“切断什么”给出了不同的回答：归一化割在没有任何先验的条件下，追求组间连接弱且两组规模平衡；图割在用户指定了部分前景和背景的条件下，把颜色模型与边界代价写成能量并求全局最小。

设把顶点集$V$分成$A$与$B$。被切断的边权之和$\operatorname{cut}(A,B)=\sum_{i\in A,j\in B}w_{i,j}$越小，两组之间的联系越弱。但直接最小化$\operatorname{cut}$会偏向切下极少数顶点：一个孤立的噪声像素只有几条弱边与其他像素相连，把它单独切出的代价最小。Shi与Malik提出\term{归一化割}（normalized cut）\scite{clust-shi2000}
\begin{equation}
  \operatorname{Ncut}(A,B)=\frac{\operatorname{cut}(A,B)}{\operatorname{assoc}(A,V)}+\frac{\operatorname{cut}(A,B)}{\operatorname{assoc}(B,V)},
  \label{eq:vision-features-ncut}
\end{equation}
其中$\operatorname{assoc}(A,V)=\sum_{i\in A,j\in V}w_{i,j}$是$A$中顶点的度之和。切下一个很小的组时，它的度之和很小，第一项接近$1$，惩罚了不平衡的划分。

算例：五个顶点，边权为$w_{1,2}=3$、$w_{3,4}=3$、$w_{2,3}=1$、$w_{4,5}=0.5$，其余为零。度依次为$3,4,4,3.5,0.5$，总和$15$。只切下顶点$5$，$\operatorname{cut}=0.5$，是所有划分中最小的，但$\operatorname{Ncut}=0.5/0.5+0.5/14.5\approx1.03$；切成$\{1,2\}$与$\{3,4,5\}$，$\operatorname{cut}=1$，$\operatorname{Ncut}=1/7+1/8\approx0.27$。归一化割选择后者，这也是直观上合理的划分：两个强连接的对子被分开，弱连接的顶点$5$跟随与它相连的顶点$4$。

精确最小化$\operatorname{Ncut}$是NP困难的组合问题。Shi与Malik把划分指示向量放松为实数向量，问题化为广义特征值问题
\[
  (\bm D-\bm W)\bm y=\lambda\bm D\bm y,
\]
其中$\bm W$为权重矩阵，$\bm D$为度构成的对角矩阵，$\bm D-\bm W$是图拉普拉斯矩阵。最小特征值$0$对应常向量，不给出划分；第二小的特征值对应的特征向量给出松弛解，再取阈值离散化为两组，递归或同时使用多个特征向量得到多组划分。卷一已经说明，这一松弛来自瑞利商的变分刻画：对满足$\bm y^{\top}\bm D\bm1=0$的实向量，$\bm y^{\top}(\bm D-\bm W)\bm y/\bm y^{\top}\bm D\bm y$的最小值就是第二小的广义特征值，连续解只是离散问题的近似，阈值化之后不保证得到最优的离散划分，可参考\BookChapterRef{chap:graph-theory}{卷一“图的表示与计算”章}。归一化割的代价主要在特征向量的计算：像素级的图有几十万个顶点，即使$\bm W$稀疏，求解也相当耗时，因此实践中常先用超像素减少顶点数。

归一化割完全自动，但“哪一块是我要的车”无法事先说明。交互式分割让用户给出少量提示，例如在车上画几笔作为前景、在路面上画几笔作为背景，再由算法完成其余像素的标记。\term{图割}（graph cut）方法用最小割求解这类问题。Boykov与Jolly把交互式分割写成二值标记的能量最小化\scite{vision-boykov-jolly2001}：为每个像素$p$选标签$\ell_p\in\{0,1\}$（背景或前景），最小化
\begin{equation}
  E(\bm\ell)=\lambda\sum_{p}D_p(\ell_p)+\sum_{\{p,q\}\in\mathcal N}B_{p,q}\,[\ell_p\neq\ell_q],
  \label{eq:vision-features-graph-cut}
\end{equation}
其中$[\cdot]$为示性函数，$\mathcal N$为相邻像素对。一元项$D_p(\ell)=-\ln\Pr(I_p\mid\ell)$衡量像素$p$的亮度或颜色与前景、背景模型的符合程度，两个模型由用户标出的像素估计。成对项$B_{p,q}$在两个相邻像素颜色相近时大、相差悬殊时小，常取$B_{p,q}\propto\exp(-(I_p-I_q)^2/(2\sigma^2))/\operatorname{dist}(p,q)$，它使分割边界倾向于落在颜色变化大的位置。用户标出的像素作为硬约束：它们的标签固定，在网络中表现为容量无穷大的终端边。

这个能量的成对项只在相邻标签不同时付出非负代价，满足卷一所说的次模条件。卷一已经证明，满足次模条件的二值成对能量可以构造为一个$s$-$t$网络，每个割的容量等于对应标记的能量加一个常数，于是最小割（用最大流算法求得）就是能量的全局最小，证明见前面提到的卷一同一章。具体构造是：每个像素一个顶点，加上源点$s$（前景）和汇点$t$（背景）；$s$到$p$的边容量为$\lambda D_p(0)$，$p$到$t$的边容量为$\lambda D_p(1)$，相邻像素之间的边容量为$B_{p,q}$。一个割把每个像素分到$s$侧或$t$侧，被切断的终端边恰好支付该像素所选标签的一元代价，被切断的邻接边恰好支付边界代价。

一个三像素链的算例说明成对项的作用。设一元代价（前景，背景）为：$p_1$为$(0.1,0.9)$，$p_2$为$(0.55,0.45)$，$p_3$为$(0.9,0.1)$；$\lambda=1$，$B_{1,2}=0.3$（颜色相近），$B_{2,3}=0.05$（对比度高）。只看一元项，$p_2$应为背景。枚举8种标记，前景、前景、背景的能量为$0.1+0.55+0.1+0.05=0.80$，前景、背景、背景为$0.1+0.45+0.1+0.3=0.95$，其余标记的能量都更高（全背景为$1.45$，全前景为$1.55$）。最优解把$p_2$标为前景：它与$p_1$颜色相近，分开要付$0.3$，而与$p_3$之间对比度高，在那里切开只需$0.05$。边界被放到了对比度最高的地方。

GrabCut在此基础上减少了用户的工作，并迭代改进颜色模型\scite{vision-rother2004}。用户只需画一个包围目标的矩形：矩形外的像素是确定的背景，矩形内的像素待定。前景与背景的颜色各用一个$K=5$个分量的全协方差高斯混合模型描述，混合模型的估计见卷四同一章。能量的一元项改为像素颜色在所属高斯分量下的负对数似然（含混合权重与协方差行列式），成对项为$\gamma\exp(-\beta\lVert\bm z_m-\bm z_n\rVert_2^2)$，其中$\beta=(2\langle\lVert\bm z_m-\bm z_n\rVert_2^2\rangle)^{-1}$由整幅图相邻像素颜色差的平均值决定，使指数项在图像的典型对比度上切换；$\gamma=50$是原文在训练图像上调出的常数。

\begin{algorithm}[GrabCut]
  \label{alg:vision-features-grabcut}
  \begin{algorithmic}[1]
    \Require 彩色图像，用户矩形（矩形外为背景$T_B$，矩形内为待定$T_U$）
    \Ensure 前景掩码
    \State 初始化：$T_B$中的像素标为背景，$T_U$中的像素标为前景，分别用两组像素初始化背景与前景的高斯混合模型
    \Repeat
      \State 对$T_U$中的每个像素，在其当前标签对应的混合模型中选取使一元代价最小的高斯分量
      \State 用分配给各分量的像素重新估计各分量的权重、均值和协方差
      \State 构造$s$-$t$网络，用最小割求$T_U$中像素的新标签
    \Until{能量不再明显下降}
    \State 用户可补画前景或背景笔画作为硬约束，再执行一次最小割
  \end{algorithmic}
\end{algorithm}

GrabCut的三步分别对分量指派、模型参数和分割标签最小化同一个能量，每一步都不增加能量，所以能量单调下降，算法收敛到局部最小。这里的局部最小只针对三组变量的交替优化：在固定颜色模型时，每次最小割仍是全局最优。它在目标与背景颜色分布差异明显时效果很好；当二者颜色重叠（例如灰色车辆停在灰色路面上），颜色模型无法区分，结果会漏掉或多出大片区域，此时需要用户补画笔画。

图割与GrabCut的交互方式在学习方法中得到了延续。\ref{sec:vision-recognition-sam}节的SAM同样接受点、框等提示并输出掩码，但它不再为每幅图重新估计颜色模型，而是由在大规模数据上训练的网络直接把提示与图像特征结合起来，能处理颜色重叠、需要形状先验的情形。图割中“一元项加成对平滑项”的能量形式，也常作为学习式分割的后处理，用来让预测边界贴合图像边缘。

表\ref{tab:vision-features-segmentation-compare}在同一街景输入下比较四类分割方法。

\begin{table*}[tbp]
  \centering
  \small
  \begin{tabularx}{\linewidth}{@{}lXXXX@{}}
    \toprule
    比较项 & 阈值（Otsu、自适应） & 区域（生长、分裂合并、分水岭） & 聚类（K均值、均值漂移、SLIC） & 图（归一化割、图割、GrabCut） \\
    \midrule
    使用的信息 & 单像素灰度 & 灰度一致性、空间邻接、梯度地形 & 颜色与位置的特征距离 & 成对相似度；图割另加用户约束与颜色模型 \\
    决策范围 & 全局或局部窗口 & 从种子或标记局部扩展 & 特征空间中的全局或局部分组 & 全图的全局优化 \\
    需要的输入 & 阈值规则 & 种子、阈值或标记 & 组数或带宽，紧致度 & 边权参数；图割需前景与背景提示 \\
    计算量 & 最低 & 低 & 中（SLIC与像素数成正比） & 归一化割高；图割中等 \\
    主要失效 & 光照不均、类别重叠 & 依赖种子与标记；过分割 & 区域不必连通；组数难定 & 归一化割偏向平衡划分；颜色重叠时图割失效 \\
    \bottomrule
  \end{tabularx}
  \caption{四类经典分割方法的比较。从左到右，决策所依据的信息逐渐增多，计算量与对用户输入的依赖也随之变化；超像素常作为图方法的预处理。}
  \label{tab:vision-features-segmentation-compare}
\end{table*}

\section{基于人工特征的识别}
\label{sec:vision-features-recognition}

有了特征，就可以回答本章开头最后一个问题：画面中有没有行人，在哪里。识别任务在这里有两种形式。检测要求输出目标的位置，通常是一组外接框；分类或检索要求判断整幅图属于哪一类、或者与图库中哪幅图最相似。经典检测采用\term{滑动窗口}（sliding window）框架：在图像金字塔的每一层上，以固定大小的窗口按一定步长扫过所有位置，对每个窗口计算特征并由分类器打分，得分超过阈值的窗口经非极大值抑制去掉重叠后作为检测结果。金字塔使固定大小的窗口能覆盖不同大小的目标，见\ref{sec:vision-transforms-pyramid}节。分类与检索则要把一幅图中数量不定的局部特征汇总成一个固定长度的向量。

这一框架下有三个相互制约的难点。第一，窗口数量巨大：一幅$640\times480$的图像在十几个尺度上可以产生数十万个窗口，而其中绝大多数是背景，每个窗口的计算必须很便宜。第二，同一类目标的外观变化大：姿态、视角、衣着与光照都会改变，特征和分类器必须容忍这些变化，同时又能与背景区分。第三，汇总局部特征时需要决定保留多少空间信息：保留越多，对目标的描述越精确，对形变和视角就越敏感。

本节的四种方法对这些难点给出了不同的回答。模板匹配直接用像素比较，只对整体亮度与对比度变化不敏感，适合外观固定的目标。Viola–Jones检测器用极其廉价的矩形特征和级联分类器，把绝大多数背景窗口在最初几步就拒绝掉，从而解决窗口数量问题；它适合外观稳定的正面人脸。HOG加线性支持向量机用梯度方向直方图容忍外观变化，可变形部件模型（DPM）进一步把目标拆成可以相对移动的部件，以应对姿态和形变。视觉词袋与VLAD面向分类与检索，把局部特征量化后汇总为直方图或残差向量，在汇总中放弃大部分空间位置，换取对视角和形变的稳健性。

\subsection{模板匹配}
\label{sec:vision-features-template-matching}

如果已知目标的外观，例如一张招牌的标准图样$t(u,v)$（大小为$m\times n$），最直接的检测方法是把它在图像上逐位置比较。平方差之和
\begin{align*}
  \operatorname{SSD}(x,y)&=\sum_{u,v}\bigl[I(x+u,y+v)-t(u,v)\bigr]^2\\
  &=\sum_{u,v}I(x+u,y+v)^2-2\sum_{u,v}I(x+u,y+v)\,t(u,v)+\sum_{u,v}t(u,v)^2
\end{align*}
越小越匹配。展开式中第三项与位置无关，中间项是图像与模板的互相关；第一项随位置变化，所以单独用互相关最大来判断匹配并不可靠，亮的区域即使与模板形状无关也会得到大的互相关。SSD本身也对光照敏感：同一招牌在阴影中整体变暗，SSD就会变大。

\term{归一化互相关}（normalized cross-correlation, NCC）在每个位置先减去窗口均值、再除以标准差：
\begin{equation}
  \operatorname{NCC}(x,y)=\frac{\sum_{u,v}\bigl(I(x+u,y+v)-\bar I_{x,y}\bigr)\bigl(t(u,v)-\bar t\bigr)}
  {\sqrt{\sum_{u,v}\bigl(I(x+u,y+v)-\bar I_{x,y}\bigr)^2\sum_{u,v}\bigl(t(u,v)-\bar t\bigr)^2}},
  \label{eq:vision-features-ncc}
\end{equation}
其中$\bar I_{x,y}$为当前窗口的均值，$\bar t$为模板均值。它就是窗口像素与模板像素之间的样本相关系数。由柯西—施瓦茨不等式，$\operatorname{NCC}\in[-1,1]$；把窗口灰度改为$aI+b$（$a>0$），减均值消去$b$，除以标准差消去$a$，NCC不变。

算例：一维模板$t=(1,2,3)$。窗口$(12,14,16)$比模板亮且对比度大一倍，去均值后为$(-2,0,2)$，与模板去均值后的$(-1,0,1)$成正比，NCC$=1$，而SSD为$11^2+12^2+13^2=434$，很大。窗口$(3,2,1)$的灰度范围与模板相同，但方向相反，NCC$=-1$，SSD只有$8$。窗口$(2,2,2)$是平坦的，标准差为零，NCC没有定义：平坦区域与任何模板都谈不上相关，实现时需要对分母加下限或直接跳过这样的窗口。

直接按定义计算，每个位置需要$mn$次乘加。分子中的互相关可以借助卷积定理用快速傅里叶变换一次算出所有位置，这在卷一已经讨论；分母中每个窗口的$\sum I$与$\sum I^2$则可以用\ref{sec:vision-features-viola-jones}节的积分图在常数时间内得到。

模板匹配的假设是目标只发生平移和仿射亮度变化。旋转几度、缩放一成、视角稍变，或者目标有部分遮挡，相关系数就会明显下降。处理旋转和尺度只能穷举一组旋转、缩放后的模板，计算量成倍增加。因此模板匹配适合工业检测、文档中的固定符号定位等外观受控的场景，或者在短时间内外观变化很小的视频跟踪。\ref{sec:vision-video-correlation-filter}节的相关滤波跟踪器在频域上高效地学习一个判别式模板，\ref{sec:vision-video-siamrpn}节的孪生网络则在学习到的特征图上做互相关，把“比较什么”交给了训练；\ref{sec:vision-video-raft}节的RAFT计算两帧所有像素特征之间的相关，也是同一运算的推广。

\subsection{Viola–Jones}
\label{sec:vision-features-viola-jones}

人脸检测曾是滑动窗口框架最典型的应用。Viola与Jones的检测器在2001年的普通处理器上达到了实时的正面人脸检测速度\scite{vision-viola2001}。它的设计完全围绕窗口数量这个难点：大多数窗口是背景，应该用尽可能少的计算把它们排除，只把计算集中在少数像人脸的窗口上。为此它组合了三个部分：可以在常数时间内计算的矩形特征、从海量候选特征中挑选少数有效特征的AdaBoost，以及逐级过滤的级联结构。

\term{类Haar特征}（Haar-like feature）是窗口内相邻矩形区域的像素和之差。两矩形特征是左右或上下两个等大矩形之差，三矩形特征是两侧之和减去中间，四矩形特征是对角两对之差。人眼区域比脸颊暗、鼻梁比两侧的眼睛亮，这样的明暗关系正好可以由这些特征刻画。在$24\times24$的基本窗口中，所有位置、大小和类型的矩形特征有18万个以上，远多于窗口的576个像素。

单个矩形的像素和本来需要逐像素相加，\term{积分图}（integral image）使它变为常数时间。积分图在位置$(x,y)$处存储原图左上方矩形内的像素和：
\begin{equation}
  ii(x,y)=\sum_{x'\leq x,\,y'\leq y}i(x',y').
  \label{eq:vision-features-integral-image}
\end{equation}
它可以一遍扫描算出：行内累加$s(x,y)=s(x,y-1)+i(x,y)$，再$ii(x,y)=ii(x-1,y)+s(x,y)$。任意矩形内的和只需四次查表：设矩形的左上角紧邻位置为$(x_0,y_0)$、右下角为$(x_1,y_1)$，则矩形和为$ii(x_1,y_1)-ii(x_0,y_1)-ii(x_1,y_0)+ii(x_0,y_0)$，最后一项是被减去两次的左上方区域。两矩形特征相邻共享两个角，只需6次查表，三矩形8次，四矩形9次。计算代价与矩形大小无关，所以检测器可以直接放大特征而不必缩放图像。

图\ref{fig:vision-features-integral-image}给出一个$4\times4$的算例。要计算第1、2行与第1、2列（下标从0开始）围成的矩形之和，在积分图上取四个格子：$ii$在第2行第2列为$14$，第0行第2列为$3$，第2行第0列为$3$，第0行第0列为$1$，矩形和为$14-3-3+1=9$，与直接相加$1+4+3+1=9$一致。

\begin{figure}[htbp]
  \centering
  % 积分图教学例：左为 4x4 原图，蓝框为第 1–2 行、第 1–2 列的矩形；右为积分图，
% 四个带符号的格子给出框内和 14-3-3+1=9。行列下标从 0 开始。宽约 68mm。
\begin{tikzpicture}[x=7mm,y=-7mm,font=\figurefont\small,
  cell/.style={draw=FigureStroke,line width=.5pt,minimum size=7mm,inner sep=0pt}]
  \fill[FigureBlueFill] (.5,.5) rectangle (2.5,2.5);
  \foreach \r/\row in {0/{1,2,0,3},1/{2,1,4,0},2/{0,3,1,2},3/{1,0,2,1}}{
    \foreach \v [count=\c from 0] in \row{\node[cell] at (\c,\r) {\v};}}
  \draw[FigureBlue,line width=1.6pt] (.5,.5) rectangle (2.5,2.5);
  \node at (1.5,-1.1) {原图 $i(x,y)$};
  \begin{scope}[xshift=40mm]
    \foreach \r/\row in {0/{1,3,3,6},1/{3,6,10,13},2/{3,9,14,19},3/{4,10,17,23}}{
      \foreach \v [count=\c from 0] in \row{\node[cell] at (\c,\r) {\v};}}
    \foreach \c/\r/\s/\f in {0/0/+/FigureBlueFill,2/0/-/FigureYellowFill,0/2/-/FigureYellowFill,2/2/+/FigureBlueFill}{
      \node[cell,fill=\f] at (\c,\r) {};
    }
    \node at (0,0) {1};\node at (2,0) {3};\node at (0,2) {3};\node at (2,2) {14};
    \foreach \c/\r/\s in {0/0/+,2/0/-,0/2/-,2/2/+}{
      \node[font=\figurefont\footnotesize,text=FigureStroke] at (\c+.32,\r-.3) {$\s$};}
    \node at (1.5,-1.1) {积分图 $ii(x,y)$};
  \end{scope}
\end{tikzpicture}

  \caption{积分图的教学例。左：$4\times4$原图，蓝框为待求和的矩形（第1、2行与第1、2列）。右：积分图，每格为原图中该格及其左上方全部像素之和；蓝底格带正号、黄底格带负号，四格相加得到矩形和$9$。}
  \label{fig:vision-features-integral-image}
\end{figure}

有了18万个以上的候选特征，下一步是挑出少数真正有用的。每个弱分类器只用一个特征：$h_j(\bm x)=1$当$p_jf_j(\bm x)<p_j\theta_j$，否则为$0$，其中$f_j$为特征值，$\theta_j$为阈值，$p_j\in\{\pm1\}$决定不等号方向。AdaBoost在每一轮为所有特征分别找出最优阈值，选出加权错误率最小的一个，再提高被它分错的样本的权重，使下一轮关注难分的样本。设第$t$轮所选弱分类器的加权错误率为$\varepsilon_t$，令$\beta_t=\varepsilon_t/(1-\varepsilon_t)$，被正确分类的样本权重乘以$\beta_t$（小于1），分错的保持不变，再重新归一化；最终的强分类器为
\begin{equation}
  H(\bm x)=\Bigl[\sum_{t=1}^{T}\alpha_th_t(\bm x)\geq\tfrac12\sum_{t=1}^{T}\alpha_t\Bigr],\qquad \alpha_t=\ln\frac1{\beta_t}.
  \label{eq:vision-features-adaboost}
\end{equation}
初始权重对$m$个负样本各取$1/(2m)$、对$l$个正样本各取$1/(2l)$，使正负两类的总权重相等。AdaBoost的指数损失解释与训练误差的界在卷四已有推导，可参考\BookChapterRef{chap:ensemble-learning}{卷四“集成学习”章}。这里它同时完成了两件事：特征选择（每一轮只选一个特征）和分类器训练。Viola与Jones报告，早期轮次选出的特征错误率在$0.1$至$0.3$之间，后期选出的在$0.4$至$0.5$之间。

单个强分类器若要达到很低的误检率，需要很多特征，每个窗口的计算仍然昂贵。\term{级联}（cascade）把检测拆成一系列由少到多的强分类器：第一级只用很少的特征，调低阈值使几乎所有人脸都能通过，同时拒绝一部分背景；通过的窗口交给第二级，依此类推，只有通过全部各级的窗口才被报告为人脸。每一级都只用前一级放过的负样本（即前一级的误检）训练，所以后面的级面对的是越来越难的背景。设第$i$级的检测率和误检率分别为$d_i,f_i$，$K$级级联的总检测率与总误检率为
\[
  D=\prod_{i=1}^{K}d_i,\qquad F=\prod_{i=1}^{K}f_i.
\]
例如10级、每级$d_i=0.99$、$f_i=0.3$，则$D=0.99^{10}\approx0.904$，$F=0.3^{10}\approx5.9\times10^{-6}$：每级的误检率并不低，但连乘后极小，代价是检测率也连乘地下降，所以每级的检测率必须非常接近1。级联还大幅降低了平均计算量：若第一级用1个特征拒绝$50\%$的窗口，第二级用10个特征拒绝剩余窗口的$80\%$，那么一个背景窗口平均只需计算$1+0.5\times10+0.1\times n_3+\cdots$个特征，其中$n_3$为第三级的特征数，后面各项的系数迅速变小（这里的拒绝比例为教学假设）。

\begin{algorithm}[级联检测器的训练]
  \label{alg:vision-features-cascade}
  \begin{algorithmic}[1]
    \Require 正样本集$P$，负样本图像集，每级目标$d_{\min},f_{\max}$，总体目标$F_{\mathrm{target}}$
    \Ensure 级联分类器
    \State $F\gets1$，从负样本图像中随机取窗口作为负样本集$N$
    \While{$F>F_{\mathrm{target}}$}
      \State 新建一级，特征数$n\gets0$
      \Repeat
        \State $n\gets n+1$，用AdaBoost在$P$与$N$上训练含$n$个特征的强分类器
        \State 在验证集上调低该级阈值，直到检测率不低于$d_{\min}$
      \Until{该级误检率不高于$f_{\max}$}
      \State $F\gets F\times(\text{该级误检率})$
      \State 用当前级联扫描负样本图像，以其误检窗口替换$N$
    \EndWhile
  \end{algorithmic}
\end{algorithm}

原文的检测器有38级、共6061个特征，前五级分别使用1、10、25、25、50个特征；训练使用4916张对齐到$24\times24$的人脸（加上水平镜像共9832张），每级使用至多一万个负窗口，负窗口来自9544张不含人脸的图像中约3.5亿个子窗口。在MIT+CMU测试集上，平均每个窗口只计算约10个特征\scite{vision-viola2001}。

Viola–Jones检测器的成功依赖于目标外观的稳定：正面、直立的人脸具有一致的明暗结构，少量矩形特征就足以描述。对侧脸、旋转的人脸，以及行人、车辆这类外观变化大的类别，矩形特征的表达能力不够；训练也很慢，每一轮AdaBoost都要在十几万个特征上搜索最优阈值。它留下的两个思想在后来的检测器中仍然可见：一是“先用廉价的计算排除大量容易的负样本”，两阶段检测器先生成少量候选框再精细分类与此相同；二是正负样本极不平衡时，要让训练集中在难分的负样本上，\ref{sec:vision-recognition-retinanet}节的RetinaNet用焦点损失在单个网络内部降低易分负样本的权重，达到了相近的目的，而不再需要级联。

\subsection{HOG+SVM与DPM}
\label{sec:vision-features-dpm}
\label{sec:vision-features-hog-svm}

行人的外观比人脸多变得多，Dalal与Triggs把\ref{sec:vision-features-hog}节的HOG描述子交给线性\term{支持向量机}（support vector machine, SVM）分类\scite{vision-dalal2005}。线性SVM学习权重$\bm w$与偏置$b$，以$\bm w^{\top}\bm\phi+b$为得分，训练时最小化$\tfrac12\lVert\bm w\rVert_2^2+C\sum_i\max(0,1-y_i(\bm w^{\top}\bm\phi_i+b))$；第一项控制间隔，第二项是合页损失，惩罚得分没有越过间隔的样本，可参考\BookChapterRef{chap:classification}{卷四“分类”章}。他们为此建立了INRIA行人数据集：从中选取1239张$64\times128$的行人图像，加上水平镜像共2478个正样本；从1218张不含行人的图像中随机取12180个窗口作为初始负样本。先训练一个初步的检测器，用它在不含行人的训练图像上穷举扫描，收集所有被误判为行人的窗口，即\term{难负例}（hard negative），把它们加入负样本集后重新训练，得到最终检测器。这一步对性能很重要：随机抽取的负样本大多是天空、路面等明显的背景，而决定分类边界的是那些像行人的背景，例如电线杆和树干。

线性SVM在HOG上的得分有一个直接的计算解释。HOG是按单元排列的$9$维直方图，块归一化之后，窗口描述子可以看成一张以块为单位的多通道特征图；$\bm w$也可以排成同样形状。于是得分$\bm w^{\top}\bm\phi$等于特征图与一个“滤波器”在该窗口处的逐元素乘积之和，在整张特征图上滑动窗口打分，就是特征图与这个滤波器的互相关。检测因此等价于在HOG特征金字塔的每一层上做一次卷积，再取阈值和非极大值抑制。

单一刚性模板描述不了行走、弯腰、骑车等多种姿态的人，也难以兼顾正面和侧面的汽车。Felzenszwalb等人的\term{可变形部件模型}（deformable part model, DPM）在这一卷积视角上加入了两种结构\scite{vision-felzenszwalb2010}。其一是部件：一个覆盖整个目标的粗分辨率\term{根滤波器}（root filter），加上若干覆盖局部（如头、手臂、车轮）的部件滤波器；部件滤波器放在特征金字塔中分辨率高一倍的层上，因而能描述更精细的局部。其二是形变：每个部件有一个相对根位置的锚点，可以偏离锚点，偏离越远付出的代价越大。这样，部件既能各自找到最匹配的位置，整体又不会散开。

设特征金字塔为$\mathcal H$，第$i$个滤波器放在位置$p_i=(x_i,y_i,l_i)$，$l_i$为金字塔层号，$\bm\phi(\mathcal H,p_i)$为该位置处滤波器覆盖的特征。一个含$n$个部件的模型对目标假设$z=(p_0,\ldots,p_n)$的得分为
\begin{equation}
  \operatorname{score}(p_0,\ldots,p_n)=\sum_{i=0}^{n}\bm F_i^{\top}\bm\phi(\mathcal H,p_i)-\sum_{i=1}^{n}\bm d_i^{\top}\bm\phi_d(dx_i,dy_i)+b,
  \label{eq:vision-features-dpm-score}
\end{equation}
其中$\bm F_0$为根滤波器，$\bm F_i$为第$i$个部件滤波器；$(dx_i,dy_i)=(x_i,y_i)-(2(x_0,y_0)+\bm v_i)$是部件相对其锚点$\bm v_i$的位移，系数$2$来自部件层的分辨率是根层的两倍；$\bm\phi_d(dx,dy)=(dx,dy,dx^2,dy^2)$为形变特征，$\bm d_i$为形变代价系数；$b$为偏置，用于让多个模型的得分可比。当$\bm d_i=(0,0,1,1)$时，形变代价就是部件位置与锚点的平方距离。

检测时，对每个根位置要选出各部件的最佳位置。由于各部件的形变代价只与根位置有关、部件之间互不相关（这是一个以根为中心的星形结构），可以对每个部件独立地求
\[
  \mathcal D_i(x,y)=\max_{dx,dy}\bigl[R_i(x+dx,y+dy)-\bm d_i^{\top}\bm\phi_d(dx,dy)\bigr],
\]
其中$R_i$为部件滤波器在部件层上的响应图。根位置的总得分就是根滤波器响应加上各部件在对应锚点处的$\mathcal D_i$值，再加偏置。直接计算$\mathcal D_i$需要对每个位置枚举所有位移；由于形变代价是可分离的二次函数，Felzenszwalb等人使用广义距离变换，在与位置数成线性的时间内算出整张$\mathcal D_i$。

一维教学算例说明形变代价的作用。设某部件在锚点附近位移$dx=-2,-1,0,1,2$处的响应为$0.1,0.6,0.5,0.9,1.3$，形变代价只取二次项$d\,dx^2$。$d=0.1$时，减去代价后为$-0.3,0.5,0.5,0.8,0.9$，最佳位移为$+2$，部件移到响应最强的位置；$d=0.3$时，结果为$-1.1,0.3,0.5,0.6,0.1$，最佳位移变为$+1$，较大的代价把部件拉向锚点。形变系数$\bm d_i$由训练学到，它规定了每个部件可以移动多远。

为了覆盖视角差异，DPM使用多个这样的模型组成混合模型，每个模型称为一个分量（如汽车的正面与侧面），检测时取各分量得分的最大值。

训练的困难在于标注只给出目标的外接框，没有给出各部件的位置，也没有给出属于哪个分量。Felzenszwalb等人把这些未标注量当作潜变量，提出\term{潜变量支持向量机}（latent SVM）。对样本$x$，分类得分取所有潜变量取值$z\in Z(x)$上的最大值，
\begin{equation}
  f_{\bm\beta}(x)=\max_{z\in Z(x)}\bm\beta^{\top}\bm\Phi(x,z),\qquad
  L(\bm\beta)=\tfrac12\lVert\bm\beta\rVert_2^2+C\sum_{i}\max\bigl(0,1-y_if_{\bm\beta}(x_i)\bigr),
  \label{eq:vision-features-latent-svm}
\end{equation}
其中$\bm\beta$把所有滤波器、形变系数和偏置拼成一个向量，$\bm\Phi(x,z)$是对应的特征，使式\eqref{eq:vision-features-dpm-score}可以写成内积。$f_{\bm\beta}$是一族关于$\bm\beta$的线性函数的最大值，因而是凸函数。对负样本（$y_i=-1$），合页损失$\max(0,1+f_{\bm\beta}(x_i))$是凸函数的最大值，仍然是凸的；对正样本，$1-f_{\bm\beta}$是凹函数，合页损失不再是凸的。这一性质称为半凸性：一旦固定每个正样本的潜变量，正样本的得分变为$\bm\beta$的线性函数，整个目标就成为凸问题。训练因此交替进行：固定$\bm\beta$，为每个正样本选出得分最高的潜变量（部件位置与分量，要求根滤波器与标注框充分重叠）；固定这些潜变量，求解凸的SVM问题；负样本则用难负例挖掘不断更新。这种坐标下降只能保证收敛到局部最优，所以初始化很重要：原文先按标注框的宽高比把正样本分成若干组，各训练一个只有根滤波器的模型，再合并为混合模型，最后在每个根滤波器能量最高的区域贪心地放置6个部件，部件滤波器由根滤波器插值到两倍分辨率初始化\scite{vision-felzenszwalb2010}。

\begin{algorithm}[潜变量SVM训练（DPM）]
  \label{alg:vision-features-latent-svm}
  \begin{algorithmic}[1]
    \Require 带外接框的正样本图像，不含目标的负样本图像，初始模型$\bm\beta$
    \Ensure 训练好的模型参数$\bm\beta$
    \For{固定的若干轮}
      \State 对每个正样本，在与标注框充分重叠的根位置中，求使$\bm\beta^{\top}\bm\Phi(x,z)$最大的$z$并固定
      \State 用当前模型扫描负样本图像，把得分越过间隔的窗口加入难负例缓存，移除容易的负例
      \State 在固定的正样本潜变量与当前负例缓存上，优化式\eqref{eq:vision-features-latent-svm}
    \EndFor
  \end{algorithmic}
\end{algorithm}

DPM所用的HOG也做了修改：先在每个单元上计算对比度敏感的18个方向和不敏感的9个方向，再加上4个反映单元周围梯度能量的分量，共31维。特征金字塔每组（尺度翻倍）的层数在训练时取5、测试时取10，更细的尺度采样对性能有明显影响。DPM在PASCAL VOC检测基准上取得了当时领先的结果，部件与形变使它能处理姿态变化和部分遮挡。它的局限在于：模型结构（分量数、部件数、部件形状）需要人工设定；特征仍是固定的HOG，表达能力有限；训练依赖精心设计的初始化与多次难负例挖掘，过程复杂。

DPM与卷积网络之间的关系后来被明确指出。Girshick等人说明，DPM的推断过程可以写成一个卷积网络：滤波器响应是卷积层，广义距离变换是一种特殊的最大池化，混合分量的取最大是另一层最大池化\scite{vision-girshick2015-dpm}。Dai等人的可变形卷积在卷积采样位置上加入由网络预测的偏移，可变形RoI池化也让各个子区域相对规则位置移动，原文把后者与DPM对照：二者都学习部件的空间形变以提高分类得分，差别在于可变形RoI池化不显式建模部件之间的空间关系，而且整个网络端到端训练\scite{vision-dai2017}。从HOG加线性SVM到DPM，再到\ref{sec:vision-recognition-detection}节的学习式检测器，“在特征金字塔上做卷积打分”的骨架一直保留，变化的是特征和形变怎样得到。

\subsection{视觉词袋与VLAD}
\label{sec:vision-features-bovw}

现在转向分类和检索：判断一幅街景是否属于“十字路口”这一场景类别，或者在一个大图库中找到拍摄同一建筑的照片。SIFT给出了每幅图上千个局部描述子，数量随图而变，分类器和索引都需要固定长度的向量。怎样把一个大小不定的描述子集合汇总成一个向量？

\term{视觉词袋}（bag of visual words, BoVW）借用文本检索中“词袋”的做法。先在训练图像的大量描述子上运行K均值，得到$K$个聚类中心$\bm c_1,\ldots,\bm c_K$，称为\term{视觉词}（visual word），全体中心组成码本。对一幅新图，把每个描述子分配给最近的视觉词，统计各视觉词出现的次数，得到$K$维直方图。这个直方图就是图像的表示：它记录了“哪些局部外观出现了多少次”，完全舍弃了它们出现在哪里。

Sivic与Zisserman把这一表示用于视频中的物体检索\scite{vision-sivic2003}，并沿用了文本检索的两项技术。一是tf-idf加权：视觉词在本图中出现得越多（词频高）、在整个库中出现的图像越少（逆文档频率高），权重越大，于是天空、路面这类到处都有的局部外观的权重被压低。二是倒排索引：为每个视觉词记录包含它的图像列表，查询时只需访问查询图中出现的视觉词对应的列表，而不必与库中每幅图逐一比较。他们还去掉了出现频率最高的一批视觉词，作用类似于文本中的停用词。

用于分类时，词袋直方图交给SVM，常用直方图交核等适合直方图的核函数。完全舍弃位置使表示对视角与形变很稳健，但也丢掉了有用的信息：天空总在上方、路面总在下方。Lazebnik等人提出\term{空间金字塔匹配}（spatial pyramid matching）\scite{vision-lazebnik2006}：把图像在第$\ell$层划分为$2^\ell\times2^\ell$个格子，在每个格子内分别统计词袋直方图，把各层各格的直方图加权拼接。对共$L$层（$\ell=0,\ldots,L$）的金字塔，第0层的权重为$1/2^L$，第$\ell\geq1$层的权重为$1/2^{L-\ell+1}$，更细的格子中的匹配说明位置也对应，因而权重更大。他们的设置为$M=200$个视觉词、$L=2$，拼接后的维数为$200\times(1+4+16)=4200$。

词袋把每个描述子压缩为一个整数（所属视觉词的编号），量化误差很大：落在同一个视觉词里的描述子，彼此的差别全部丢失。要减小这一误差，就要用很大的码本，例如数十万个视觉词，而码本越大，K均值越难训练，直方图也越稀疏。Jégou等人提出\term{局部聚合描述子向量}（vector of locally aggregated descriptors, VLAD），用小码本保留更多信息\scite{vision-jegou2010}：不只统计每个视觉词被分到多少个描述子，而是累加这些描述子与中心的残差，
\begin{equation}
  \bm v_k=\sum_{\bm x:\,\operatorname{NN}(\bm x)=\bm c_k}(\bm x-\bm c_k),\qquad k=1,\ldots,K,
  \label{eq:vision-features-vlad}
\end{equation}
其中$\operatorname{NN}(\bm x)$是离$\bm x$最近的中心。把$\bm v_1,\ldots,\bm v_K$拼接成$Kd$维向量（$d$为描述子维数，SIFT为128），再做L2归一化。残差记录了落入同一个视觉词的描述子在该词内部的分布偏向，使较小的码本（几十个中心）就能得到有区分力的表示；随后的降维与压缩编码又把它压到很短，适合大规模检索。

算例：二维描述子、两个中心$\bm c_1=(1,1)$、$\bm c_2=(1,3)$。图甲的四个描述子为$(1,0),(2,1),(0,3),(1,4)$，前两个最近的中心是$\bm c_1$，后两个是$\bm c_2$；图乙的四个描述子为$(2,1),(0.5,1),(2,3),(1,3.5)$，同样是前两个属于$\bm c_1$、后两个属于$\bm c_2$。两幅图的词袋直方图都是$(2,2)$，余弦相似度为$1$，无法区分。VLAD中，图甲$\bm v_1=(0,-1)+(1,0)=(1,-1)$，$\bm v_2=(-1,0)+(0,1)=(-1,1)$，拼接并归一化为$(0.5,-0.5,-0.5,0.5)$；图乙$\bm v_1=(1,0)+(-0.5,0)=(0.5,0)$，$\bm v_2=(1,0)+(0,0.5)=(1,0.5)$，归一化为$(0.41,0,0.82,0.41)$。两个VLAD向量的内积为$0.204-0.408+0.204=0$，表示被清楚地分开了。

词袋与VLAD都分两步：先由无监督的K均值学码本，再固定码本做编码。码本只反映描述子的分布，并不针对检索或分类目标优化；底层的SIFT也是固定的。Arandjelović等人的NetVLAD把VLAD改写为可微的网络层\scite{vision-arandjelovic2016}：把硬分配换成软分配$\bar a_k(\bm x)=\exp(\bm w_k^{\top}\bm x+b_k)/\sum_{k'}\exp(\bm w_{k'}^{\top}\bm x+b_{k'})$，残差按$\bar a_k$加权累加，中心$\bm c_k$与分配参数$\bm w_k,b_k$都作为参数学习；输入换成卷积网络最后一个卷积层的局部特征，整个网络用地点识别的弱监督端到端训练。原文以VGG的conv5特征（512维）、$K=64$为设置，聚合后按中心逐块L2归一化（块内归一化），再整体L2归一化，并可用带白化的主成分分析降到4096维。\ref{sec:vision-recognition-gem}节的GeM则走得更远，直接对卷积特征做广义均值池化，不再使用码本。从词袋到NetVLAD与GeM，变化的是局部特征和汇总方式由人工设计变为由检索目标学习，“先提局部特征、再汇总成全局向量”的结构保持不变。

表\ref{tab:vision-features-recognition-compare}比较本节的四种方法。

\begin{table*}[tbp]
  \centering
  \small
  \begin{tabularx}{\linewidth}{@{}lXXXX@{}}
    \toprule
    比较项 & 模板匹配 & Viola–Jones & HOG+SVM与DPM & 视觉词袋与VLAD \\
    \midrule
    任务 & 定位已知外观的目标 & 检测外观稳定的类别（正面人脸） & 检测外观多变的类别（行人、车辆） & 图像分类与检索 \\
    特征 & 像素 & 类Haar矩形特征 & HOG；DPM另有部件与形变 & 量化后的局部描述子（如SIFT） \\
    分类或比较 & NCC & AdaBoost强分类器的级联 & 线性SVM；潜变量SVM & 直方图核SVM；内积与倒排索引 \\
    处理效率的方式 & FFT与积分图 & 积分图、级联提前拒绝 & 特征金字塔上的卷积、广义距离变换 & 倒排索引、短向量编码 \\
    对外观变化 & 只容忍亮度与对比度 & 容忍小的光照与位移 & 容忍光照、小形变；DPM容忍姿态 & 舍弃位置，容忍视角与形变 \\
    主要失效 & 旋转、缩放、遮挡 & 侧脸、旋转、外观多变的类别 & 大姿态变化、严重遮挡、小目标 & 空间布局信息丢失；码本不针对任务 \\
    \bottomrule
  \end{tabularx}
  \caption{基于人工特征的识别方法的比较。前三种按滑动窗口框架检测目标，第四种把局部特征汇总为全局向量。每一种都可以看出与后来学习方法的对应：孪生网络的相关、两阶段检测与焦点损失、可变形卷积、NetVLAD与GeM。}
  \label{tab:vision-features-recognition-compare}
\end{table*}

\section{本章小结}
\label{sec:vision-features-summary}

本章开头提出了四个问题：车道线在哪里，同一块招牌在另一幅照片中位于何处，路面与天空各占哪些像素，画面中有没有行人。它们都要求从像素值中提取比单个像素更稳定的量，再把这些量组织成几何对象、对应关系、区域或检测结果。

第一个问题由边缘与直线的检测回答。梯度给出灰度变化的强度与方向，二阶导数的零交叉给出不依赖阈值的位置；Canny从检测、定位与单一响应三条准则出发，说明平滑宽度只能在信噪比与定位精度之间取舍，乘积由算子形状决定，并以非极大值抑制和滞后阈值得到细而连贯的边缘。Hough变换让边缘点在参数空间中投票，把断裂的车道线重新组织为直线。第二个问题由局部特征回答。结构张量的两个特征值区分了平坦、边缘与角点；SIFT在尺度空间中寻找归一化响应的极值，使检测的尺度随图像缩放而变化，再在按尺度和方向归一化的邻域中计算梯度方向直方图，实现跨尺度、跨旋转的匹配。第三个问题由分割和形状描述回答。阈值、区域、聚类与图四类方法依次使用更多的信息：Otsu的类间方差准则可以从全方差分解严格导出，分水岭用梯度地形放置闭合边界，SLIC以局部K均值得到超像素，归一化割与图割在图上做全局优化，后者在满足次模条件时可以用最小割求得精确最优。链码、傅里叶描述子、紧致度、欧拉数与Hu矩再把区域概括为对平移、旋转、缩放不变的数字。第四个问题由基于人工特征的识别回答：积分图与级联使Viola–Jones能实时检测人脸，HOG与线性SVM在特征金字塔上做卷积打分，DPM加入部件与形变并用潜变量SVM训练，视觉词袋与VLAD把局部特征汇总为可检索的全局向量。

这些方法共同留下了几个尚未解决的问题。它们的不变性都是设计者事先选定的：SIFT对相似变换不变，却不能处理大视角；HOG容忍小形变，却要靠DPM的人工部件结构才能应对姿态。它们只使用低层线索，无法区分物体轮廓与阴影边界，也无法把颜色相同的车顶与路面分开，因为“什么是有意义的边界或区域”取决于语义。它们的参数（尺度、阈值、码本大小、部件数）需要针对数据调整，并且各个模块分别设计、分别优化，前一步的错误无法由后一步纠正。

后续各章正是在这些地方作出改变。\ref{chap:vision-video-foundations}章把结构张量、金字塔与块匹配用于估计视频中的运动。\ref{chap:vision-image-video-models}章的卷积网络与视觉Transformer把滤波器、汇总方式与归一化都交给学习，并在大规模数据上得到语义相关的特征。\ref{chap:vision-image-recognition}章再在这些特征之上完成检测、分割与检索，那里会看到本章的许多结构以可学习的形式重新出现：特征金字塔对应尺度空间，可变形卷积对应DPM的部件形变，NetVLAD与GeM对应词袋汇总，SAM的提示式分割对应图割的交互方式。

\paragraph{扩展阅读}
局部特征方向，SIFT之后的工作主要在速度与不变性之间重新取舍。Bay等人的SURF用积分图上的盒子滤波近似Hessian矩阵的行列式与Haar小波响应，在重复率接近SIFT的同时明显缩短了检测与描述的时间；Rosten与Drummond的FAST只比较候选点周围一圈像素与中心的明暗，并用决策树学习比较的顺序，适合实时跟踪；Rublee等人的ORB在FAST角点上加入Harris响应排序与灰度质心方向，配以按主方向旋转的二值描述子，用汉明距离代替SIFT的欧氏距离完成匹配\scite{vision-bay2008surf,vision-rosten2006fast,vision-rublee2011orb}Matas等人的最大稳定极值区域（maximally stable extremal regions, MSER）检测的不是点，而是阈值扫过整个灰度范围时面积变化最慢的连通区域，可以看作把\ref{sec:vision-features-otsu}节的阈值分割改造成一种对仿射变换较稳健的区域检测器。要判断这些检测器与描述子在视角、模糊和光照变化下谁更可靠，可以读Mikolajczyk与Schmid的评测，它给出了比较描述子所用的图像序列和查全率—误检率协议\scite{vision-matas2004mser,vision-mikolajczyk2005descriptors}汇总方式上，Fisher向量与VLAD一样累加局部描述子相对中心的偏差，但用高斯混合模型的软分配代替K均值的硬分配，并同时记录一阶与二阶统计量，VLAD可以看作它的简化；Perronnin等人加入幂归一化与L2归一化后，它成为深度特征普及之前大规模图像分类的代表表示\scite{vision-perronnin2010fisher}

分割方向，本章没有展开以曲线演化为核心的变分方法。Kass等人的主动轮廓（Snake）把轮廓写成一条同时受内部平滑能量和图像边缘能量约束的参数曲线，从用户给出的初始位置出发向边缘收敛；Chan与Vese用水平集表示轮廓，以轮廓内外的灰度均值差代替梯度驱动曲线演化，能分割边缘模糊的目标，轮廓的拓扑也可以在演化中自然分裂或合并\scite{vision-kass1988snakes,vision-chan2001activecontours}它们与\ref{sec:vision-features-graph-segmentation}节的图割同样最小化“数据项加平滑项”形式的能量，区别在于沿连续曲线做梯度下降，一般只能得到局部最优。Felzenszwalb与Huttenlocher的图分割按边权从小到大合并区域，比较两区域之间的最小边权与各自内部差异来决定是否合并，复杂度为$O(m\log m)$（$m$为边数），常与SLIC一样用来产生过分割；Uijlings等人的选择性搜索以这种过分割为起点，按颜色、纹理、大小与吻合程度等多种相似度逐层合并，得到与类别无关的候选框，曾是R-CNN等检测器的候选区域来源\scite{vision-felzenszwalb2004graph,vision-uijlings2013selective}评价边缘检测与分割时常用伯克利分割数据集，Arbeláez等人在其上建立了从轮廓检测到层次分割的完整流程和评价指标\scite{vision-arbelaez2011contour}若想系统回顾本章各主题，Szeliski的教材覆盖了特征检测与匹配、分割和识别，并说明了它们与学习方法的联系\scite{szeliski2022computer}
